% !TeX root = ../../Book4.tex
% 纲要标题：## 第1章　范畴、函子与自然变换
% 纲要位置：计划纲要.md，第 2449 行。

\chapter{范畴、函子与自然变换}
\label{b4:ch:01}
\BookChapterNumber{b4:ch:01}

\begin{chapterabstract}
本章从群、偏序与模等例子建立范畴、函子和自然变换，讨论自然同构与依赖选择的对象同构有何区别，并证明范畴等价的判据及若干新范畴的构造。将态射的复合线性化以后，有限对象的模图表还可以统一写成范畴代数上的左模。
\end{chapterabstract}

\begin{learninggoals}
\begin{itemize}
\item 在具体例子中核对范畴公理，区分局部小范畴与小范畴。
\item 准确写出协变和反变构造的态射方向，并检验自然性方块。
\item 计算自然变换的复合，证明自然同构可以逐分量判断。
\item 区分全、忠实和本质满，在说明选择条件后构造拟逆。
\item 使用函子、切片与逗号范畴记录图表及指定映射。
\item 辨认已有模同构中的自然部分与依赖选择的部分。
\item 构造范畴代数，并在模图表与范畴代数的左模之间恢复对象和态射。
\end{itemize}
\end{learninggoals}

有限生成交换群都可以分解为挠子群与一个自由群的直和。例如，\(\mathbb Z\oplus\mathbb Z/3\mathbb Z\) 的挠子群是 \(0\oplus\mathbb Z/3\mathbb Z\)，而 \(\mathbb Z\oplus0\) 是一个自由补群。但是，分别给每个群选定补群之后，群同态是否会把所选补群送入另一个群的所选补群？能否让所有选择都与群同态相容？分类定理本身没有回答这个问题。为研究这种相容性，我们要同时考虑对象、对象之间的映射，以及构造对这些映射的作用；范畴、函子和自然变换将把这些要求写成可以核验的等式。

\ProseParagraphBreak
主要例子需要集合与映射、群同态、环和模的基本知识。第二卷第5章曾围绕模论介绍最低限度的范畴语言；本章进一步讨论大小条件、自然变换的复合和新范畴的构造。关于自然同构的讨论通过第二卷第1.2节的双对偶、第6.2节的张量同构及第4.3节的 PID 模结构，比较具体同构对映射的相容性，并回答开头的分裂问题。范畴代数则把一个模图表的多个分量合成一个模，使态射的复合成为代数乘法。

\ProseParagraphBreak
参考阅读可从 Weibel\cite[Sections A.1--A.3, pp. 417--424]{Weibel1994HomologicalAlgebra}开始：A.1 给出范畴与大小的基本例子，A.2 讨论函子及代表选择，A.3 介绍自然变换、等价与函子范畴。本章对等价判据给出完整证明，并把一般大范畴中所需的选择数据单独写入陈述。

% !TeX root = ../../Book4.tex
% 纲要标题：### 1.1 范畴的定义与例子
% 纲要位置：计划纲要.md，第 2451 行。

\section{范畴的定义与例子}
\label{b4:sec:0101}

群、模与向量空间都通过各自的同态相互比较。把恒等映射和复合的规律单独提取出来，便能在不同代数结构中讨论同一种泛性质。

% 纲要内容：
%
% * 对象、态射及局部小性
% * 模范畴、偏序范畴及单对象范畴
% * 集合大小与工作约定
%

\subsection{对象、态射及局部小性}
\label{b4:subsec:0101:01}

研究向量空间时，线性映射不仅用来判断两个空间是否同构，还用来构造核、商空间、对偶和张量积。换成群或模，映射的具体定义改变了，恒等映射与映射复合却仍满足同样的规律。范畴先保留这些规律；此后讨论的泛性质只依赖对象之间允许哪些映射，以及这些映射怎样复合。

\begin{definition}\label{b4:def:category}
一组数据 \(\mathcal C\) 包含一个对象类 \(\operatorname{Ob}\mathcal C\)，以及每对对象 \(X,Y\) 之间的一个态射类，记为 \(\operatorname{Hom}_{\mathcal C}(X,Y)\)，也可以记为 \(\operatorname{Mor}_{\mathcal C}(X,Y)\)。每个对象还带有恒等态射 \(1_X\in\operatorname{Hom}_{\mathcal C}(X,X)\)，并给定复合
\[
\operatorname{Hom}_{\mathcal C}(Y,Z)\times\operatorname{Hom}_{\mathcal C}(X,Y)\longrightarrow\operatorname{Hom}_{\mathcal C}(X,Z),
\qquad (g,f)\longmapsto g\circ f.
\]
每个态射带有指定的起点和终点。若复合满足
\[
h\circ(g\circ f)=(h\circ g)\circ f,\qquad
1_Y\circ f=f=f\circ1_X
\]
对所有可复合态射成立，则称这些数据组成一个\term[fanchou]{范畴}[category]。若每个 \(\operatorname{Hom}_{\mathcal C}(X,Y)\) 都是集合，称 \(\mathcal C\) 为\term[jubuxiaofanchou]{局部小范畴}[locally small category]。
\end{definition}
\symidx{\(\operatorname{Hom}_{\mathcal C}(X,Y)\)：范畴中的态射类}
\symidx{\(\operatorname{Mor}_{\mathcal C}(X,Y)\)：范畴中的态射类的别称记号}

复合 \(gf\) 始终表示先作 \(f\)、再作 \(g\)。结合律只在各箭头首尾相接时使用；它不允许复合两个终点、起点不匹配的态射。本卷以后所说的范畴默认局部小，涉及更大范围时另作说明。空的对象类也允许，它给出空范畴。

\ProseParagraphBreak
定义没有要求态射一定是底集合之间的函数。许多代数范畴确实如此，但偏序与单对象范畴会说明，箭头也可以记录一个关系或一个代数元素。范畴公理讨论的是箭头的复合，不能在尚未给定底集合时直接取“一个对象的元素”。

\begin{definition}\label{b4:def:categorical-isomorphism}
设 \(\mathcal C\) 为范畴，\(f:X\rightarrow Y\) 为其中态射。若存在态射 \(g:Y\rightarrow X\)，使
\(gf=1_X\)、\(fg=1_Y\)，则称 \(f\) 为同构，并记 \(X\cong Y\)。
\end{definition}

逆态射唯一：若 \(g,h\) 都是 \(f\) 的逆，则
\(g=g(fh)=(gf)h=h\)。恒等态射是同构，同构的复合仍为同构，且
\((gf)^{-1}=f^{-1}g^{-1}\)。这使“对象同构”成为对象之间的等价关系，但不使所有同构对象成为同一个对象。

\begin{definition}\label{b4:def:monomorphism-epimorphism}
设 \(\mathcal C\) 为范畴，\(f:X\rightarrow Y\) 为态射。
若对每个对象 \(T\) 及态射 \(u,v:T\rightarrow X\)，等式 \(fu=fv\) 都推出 \(u=v\)，称 \(f\) 为\term[dantaishe]{单态射}[monomorphism]。
若对每个对象 \(T\) 及态射 \(u,v:Y\rightarrow T\)，等式 \(uf=vf\) 都推出 \(u=v\)，称 \(f\) 为\term[mantaishe]{满态射}[epimorphism]。
\end{definition}

单态射是可在等式左侧消去的箭头，满态射是可在右侧消去的箭头。它们不是关于对象元素的定义。同构同时是单态射与满态射，反过来在一般范畴中并不成立。

\subsection{模范畴、偏序范畴及单对象范畴}
\label{b4:subsec:0101:02}

集合及函数组成范畴 \(\mathbf{Set}\)，群及群同态组成 \(\mathbf{Grp}\)，交换群组成 \(\mathbf{Ab}\)。含幺环及保幺环同态组成 \(\mathbf{Ring}\)，其中交换环组成 \(\mathbf{CRing}\)；这里允许零环。恒等映射保持相应结构，保持结构的映射之复合仍保持结构，结合律来自函数复合，因此它们都满足范畴公理。

\ProseParagraphBreak
固定含幺环 \(R\)，幺性左 \(R\) 模及 \(R\) 线性映射组成
\(R\text{-}\mathbf{Mod}\)。任意固定的两个模 \(M,N\) 之间的同态是函数集合 \(N^M\) 的一个子集，所以这个范畴局部小。域 \(k\) 上的模范畴也记为 \(\mathbf{Vect}_k\)；只保留有限维对象，得到 \(\mathbf{Vect}_k^{\mathrm{fd}}\)。

\begin{proposition}\label{b4:prop:module-monos-epis}
设 \(R\) 为含幺环，\(f:M\rightarrow N\) 为幺性左 \(R\) 模同态。在 \(R\text{-}\mathbf{Mod}\) 中，\(f\) 为单态射当且仅当它作为函数单射；\(f\) 为满态射当且仅当它作为函数满射。因此在这个范畴中，同时为单态射与满态射的态射就是同构。
\end{proposition}
\begin{proof}
单射或满射函数分别具有相应的消去性质，所以两个充分性方向成立。若 \(f\) 不是单射，取 \(0\ne m\in\ker f\)。映射
\(a:R\rightarrow M\)、\(a(r)=rm\)，与零映射不同，却满足 \(fa=0\)，所以 \(f\) 不是单态射。

若 \(f\) 不是满射，则商模 \(Q=N/f(M)\ne0\)。商映射
\(q:N\rightarrow Q\) 与零映射不同，而 \(qf=0\)，所以 \(f\) 不是满态射。最后，双射模同态的集合逆保持加法与标量乘法：把所需等式施以 \(f\)，两边相等，再用单射性即可。因此它是模同构。
\end{proof}

不能把最后一句原样用于环范畴。保幺包含 \(\mathbb Z\hookrightarrow\mathbb Q\) 在 \(\mathbf{CRing}\) 中既单又满：单性来自底函数单射；若
\(a,b:\mathbb Q\rightarrow S\) 在 \(\mathbb Z\) 上相同，则每个分数的像都由
\[
a(m/n)=a(m)\,a(n)^{-1}
\]
唯一确定，故 \(a=b\)。然而这个包含不是环同构，因为它不是满射。这项例子也说明，满态射要求检测的目标及允许的映射必须属于指定范畴。

\begin{example}\label{b4:ex:poset-category}
设 \((P,\le)\) 为偏序集。以 \(P\) 的元素为对象，规定当 \(x\le y\) 时有唯一箭头 \(x\rightarrow y\)，否则没有箭头。反身性给出恒等箭头，传递性给出复合；任意两个相同起终点的箭头都相等，所以结合律自动成立。同构 \(x\cong y\) 要求同时有 \(x\le y\)、\(y\le x\)，故只在 \(x=y\) 时成立。每个箭头都单且满，因为所有需要比较的平行箭头本来就至多一个。
\end{example}

若只要求关系反身、传递，不要求反对称，仍得到范畴，称为预序范畴；此时不同对象可以同构。比如两个元素各自都与另一个元素满足关系，就给出两个不同而同构的对象。一个集合还可以看作只有恒等态射的离散范畴。

\begin{example}\label{b4:ex:one-object-category}
设 \(M\) 为含单位元 \(1\) 的幺半群。只取一个对象 \(*\)，规定
\(\operatorname{Hom}_{\mathcal BM}(*,*)=M\)，并令 \(b\circ a=ba\)，就得到单对象范畴 \(\mathcal BM\)。反过来，一个局部小单对象范畴的全部自态射，在复合下组成幺半群。同构态射恰好对应幺半群的可逆元；当 \(M=G\) 为群时，每个箭头都可逆。
\end{example}

这三个例子展示了相同公理的不同用法：模范畴保留所有线性映射，偏序范畴记录何时存在一个比较，单对象范畴则把整个乘法放进自态射集合中。后文的函子必须分别尊重线性映射的复合、次序关系或幺半群乘法。

\subsection{集合大小的约定}
\label{b4:subsec:0101:03}

\begin{definition}\label{b4:def:small-category}
设 \(\mathcal C\) 为局部小范畴。若 \(\operatorname{Ob}\mathcal C\) 为集合，称 \(\mathcal C\) 为\term[xiaofanchou]{小范畴}[small category]；此时全部态射也是集合，因为它们是集合指标的一族 Hom 集的并。若对象类不为集合，称 \(\mathcal C\) 为大范畴。
\end{definition}

“局部小”只控制每一对对象之间的态射，“小”还控制对象总数。偏序集和幺半群构成的小范畴满足后一条件；全部集合、全部群或全部 \(R\) 模组成的范畴通常只局部小。若试图把所有集合收进一个集合，便会遇到集合论悖论，所以不能省略这项区别。

\ProseParagraphBreak
本卷使用通常的集合与类语言，并允许集合指标的选择公理。写 \(\mathbf{Set}\) 时是指所讨论范围中的集合范畴，不把它的对象类当作一个可任意取幂集的集合。形成函子范畴 \(\operatorname{Fun}(\mathcal C,\mathcal D)\) 时，通常要求源范畴 \(\mathcal C\) 小、目标 \(\mathcal D\) 局部小；这样两函子之间的自然变换构成集合，具体证明见\cref{b4:prop:functor-category-locally-small}。

\ProseParagraphBreak
另一个需要大小约定的地方是“为每个对象选一个同构代表”。若要选择的对象组成集合，通常的选择公理可以使用；真类指标的选择并不自动由此得到。涉及大范畴等价时，本卷会给出具体的拟逆，或把所需代表及同构作为已选数据。若采用固定宇宙的写法，也必须在更大的工作范围中进行这些集合选择；这是一种大小安排，不能代替对选择步骤的说明。

\ProseParagraphBreak
这几条约定已经足够本编使用。此后每个“任意族”“所有函子”或“逐对象选择”都有指定范围；不需要为了学习范畴论先建立另一套集合论，却也不能把真类当作普通集合计算。

% !TeX root = ../../Book4.tex
% 纲要标题：### 1.2 函子与自然变换
% 纲要位置：计划纲要.md，第 2457 行。

\section{函子与自然变换}
\label{b4:sec:0102}

一种构造若要适用于全部同态，就必须同时说明它怎样改变对象和映射。比较两种这样的构造时，所选分量还须与每个原同态交换。

% 纲要内容：
%
% * 协变、反变及自然性方块
% * 函子复合与自然变换复合
% * 忘却结构和附加结构的例子
% * 生成态射上的自然性检验；逐对象同构但不自然同构的两对象反例
%

\subsection{协变函子、反变函子与自然性方块}
\label{b4:subsec:0102:01}

把一个模忘掉标量作用后变成交换群，是一种对对象的操作；把模同态也看成群同态，才说明这种操作怎样与映射相容。函子同时规定这两部分。只列出每个对象要变成什么，还没有给出函子。

\begin{definition}\label{b4:def:functor}
设 \(\mathcal C,\mathcal D\) 为范畴。给每个 \(X\in\mathcal C\) 指定对象 \(F(X)\in\mathcal D\)，给每个态射 \(f:X\rightarrow Y\) 指定态射
\(F(f):F(X)\rightarrow F(Y)\)。若
\[
F(1_X)=1_{F(X)},\qquad F(g\circ f)=F(g)\circ F(f)
\]
对所有对象及可复合态射成立，称这组指定为从 \(\mathcal C\) 到 \(\mathcal D\) 的\term[xiebianhanzi]{协变函子}[covariant functor]，记为 \(F:\mathcal C\rightarrow\mathcal D\)。本书简称“函子”时指协变函子。
\end{definition}

\begin{definition}\label{b4:def:contravariant-functor}
设 \(\mathcal C,\mathcal D\) 为范畴。给每个对象 \(X\in\mathcal C\) 指定对象 \(F(X)\in\mathcal D\)，给每个 \(\mathcal C\) 中的态射 \(f:X\rightarrow Y\) 指定反向态射
\(F(f):F(Y)\rightarrow F(X)\)。若
\[
F(1_X)=1_{F(X)},\qquad F(g\circ f)=F(f)\circ F(g),
\]
对所有对象及可复合态射成立，称这组指定为从 \(\mathcal C\) 到 \(\mathcal D\) 的\term[fanbianhanzi]{反变函子}[contravariant functor]。
\end{definition}

写出一个具体函子时，要同时说明它把对象送到什么对象、把态射送到什么态射。下面以两个 Hom 函子为例，分别写出这两种对应。

\ProseParagraphBreak
固定局部小范畴 \(\mathcal C\) 中的对象 \(A\)。从 \(A\) 到每个对象 \(X\) 的全体态射组成集合 \(\operatorname{Hom}_{\mathcal C}(A,X)\)。先写出协变 Hom 函子在对象上的对应：
\[
\begin{aligned}
\operatorname{Hom}_{\mathcal C}(A,-):\mathcal C&\longrightarrow\mathbf{Set},\\
X&\longmapsto\operatorname{Hom}_{\mathcal C}(A,X).
\end{aligned}
\]

再给定态射 \(f\in\operatorname{Hom}_{\mathcal C}(X,Y)\)。它在这个函子作用下对应到下面的集合映射：
\[
\begin{aligned}
\operatorname{Hom}_{\mathcal C}(A,-)(f):
\operatorname{Hom}_{\mathcal C}(A,X)&\longrightarrow
\operatorname{Hom}_{\mathcal C}(A,Y),\\
u&\longmapsto f\circ u.
\end{aligned}
\]

因此，固定两个对象 \(X,Y\)，该函子在它们之间的态射集上的作用可以写为
\[
\begin{aligned}
\operatorname{Hom}_{\mathcal C}(A,-):
\operatorname{Hom}_{\mathcal C}(X,Y)&\longrightarrow
\operatorname{Hom}_{\mathbf{Set}}\bigl(
\operatorname{Hom}_{\mathcal C}(A,X),
\operatorname{Hom}_{\mathcal C}(A,Y)\bigr),\\
f&\longmapsto\bigl(u\longmapsto f\circ u\bigr).
\end{aligned}
\]
这里右端是两个集合之间的全体映射；左端的每个态射 \(f\) 被送到其中一个映射，而 \(u\) 是这个映射的输入。

\ProseParagraphBreak
符号中的横线 \(-\) 表示随对象变化的位置，\(A\) 始终固定。这里能以 \(\mathbf{Set}\) 为目标范畴，正是因为 \(\mathcal C\) 局部小，每个 Hom 类都是集合。

\ProseParagraphBreak
另一种做法是取从 \(X\) 到 \(A\) 的全体态射。相应的反变 Hom 函子在对象上的对应为
\[
\begin{aligned}
\operatorname{Hom}_{\mathcal C}(-,A):\mathcal C&\longrightarrow\mathbf{Set}
\quad\text{（反变）},\\
X&\longmapsto\operatorname{Hom}_{\mathcal C}(X,A).
\end{aligned}
\]

给定态射 \(f\in\operatorname{Hom}_{\mathcal C}(X,Y)\)，它对应到的集合映射为
\[
\begin{aligned}
\operatorname{Hom}_{\mathcal C}(-,A)(f):
\operatorname{Hom}_{\mathcal C}(Y,A)&\longrightarrow
\operatorname{Hom}_{\mathcal C}(X,A),\\
v&\longmapsto v\circ f.
\end{aligned}
\]
它把从 \(Y\) 到 \(A\) 的态射 \(v\) 送到从 \(X\) 到 \(A\) 的态射 \(v\circ f\)，所以映射的方向与 \(f\) 相反。

\ProseParagraphBreak
于是，对固定的 \(X,Y\)，该函子在态射集上的作用为
\[
\begin{aligned}
\operatorname{Hom}_{\mathcal C}(-,A):
\operatorname{Hom}_{\mathcal C}(X,Y)&\longrightarrow
\operatorname{Hom}_{\mathbf{Set}}\bigl(
\operatorname{Hom}_{\mathcal C}(Y,A),
\operatorname{Hom}_{\mathcal C}(X,A)\bigr),\\
f&\longmapsto\bigl(v\longmapsto v\circ f\bigr).
\end{aligned}
\]

恒等律分别来自 \(1_X\circ u=u\)、\(v\circ1_X=v\)，复合规律分别来自
\[
(g\circ f)\circ u=g\circ(f\circ u),\qquad
v\circ(g\circ f)=(v\circ g)\circ f.
\]
这类 Hom 函子的系统作用将在第3章讨论。

\ProseParagraphBreak
域 \(k\) 上的代数对偶也是反变函子：\(V^*=\operatorname{Hom}_k(V,k)\)，而
\(f^*(\lambda)=\lambda f\)。若 \(V\xrightarrow{f}W\xrightarrow{g}Z\)，则对 \(\mu\in Z^*\)，
\((gf)^*(\mu)=\mu gf=f^*(g^*(\mu))\)。第二卷的转置矩阵规则正是这个复合次序在坐标中的表达。
\symidx{\(V^*\)：向量空间的代数对偶}

\begin{proposition}\label{b4:prop:functors-preserve-isomorphisms}
设 \(\mathcal C,\mathcal D\) 为范畴，\(F:\mathcal C\rightarrow\mathcal D\) 为函子。若 \(f:X\rightarrow Y\) 是 \(\mathcal C\) 中的同构，则 \(F(f)\) 是 \(\mathcal D\) 中的同构，且
\(F(f^{-1})=F(f)^{-1}\)。反变函子也保留同构。
\end{proposition}
\begin{proof}
对等式 \(f^{-1}f=1_X\)、\(ff^{-1}=1_Y\) 应用协变函子，得到 \(F(f^{-1})\) 的左右逆等式。反变情形把乘积次序调换，仍然得到这两个逆等式。
\end{proof}

\begin{definition}\label{b4:def:natural-transformation}
设 \(\mathcal C,\mathcal D\) 为范畴，\(F,G:\mathcal C\rightarrow\mathcal D\) 为两个函子。对每个对象 \(X\in\mathcal C\)，给定 \(\mathcal D\) 中的态射
\(\eta_X:F(X)\rightarrow G(X)\)。若对每个 \(f:X\rightarrow Y\)，都有
\[
G(f)\eta_X=\eta_YF(f),
\]
则称这族态射为从 \(F\) 到 \(G\) 的\term[ziranbianhuan]{自然变换}[natural transformation]，记为 \(\eta:F\Rightarrow G\)，称 \(\eta_X\) 为在 \(X\) 处的分量。
\end{definition}
\symidx{\(\eta:F\Rightarrow G\)：函子之间的自然变换}

等式可以画为自然性方块：
\[
\begin{tikzcd}
F(X)\arrow[r,"\eta_X"]\arrow[d,"F(f)"']&
G(X)\arrow[d,"G(f)"]\\
F(Y)\arrow[r,"\eta_Y"']&G(Y).
\end{tikzcd}
\]
它比较两种次序：先用分量映射再沿 \(f\) 变化，或者先沿 \(f\) 变化再用分量映射。自然性要求对源范畴中的全部态射成立，但证明时不一定要逐一检验。若给定的一族态射生成源范畴，即每个态射都是这族态射的有限复合或恒等态射，则只需在生成态射上检验。事实上，恒等态射处的自然性自动成立；若分量族对 \(f:X\rightarrow Y\)、\(g:Y\rightarrow Z\) 都满足自然性，则
\[
\begin{aligned}
G(gf)\eta_X
&=G(g)G(f)\eta_X
=G(g)\eta_YF(f)\\
&=\eta_ZF(g)F(f)
=\eta_ZF(gf).
\end{aligned}
\]
因此自然性也对它们的复合成立。任意抽选若干态射，或只检验同构，通常不足以推出自然性。

\ProseParagraphBreak
两个反变函子之间的自然变换使用同样的要求，但两个函子所给的态射方向都与原箭头 \(f:X\rightarrow Y\) 相反，因此等式为
\(G(f)\eta_Y=\eta_XF(f)\)。普通自然变换要求两个函子有相同的源范畴和目标范畴。\secref{b4:sec:0103}将把反变函子表述为从对偶范畴 \(\mathcal C^{\mathrm{op}}\) 到 \(\mathcal D\) 的协变函子；要把它与 \(\mathcal C\rightarrow\mathcal D\) 比较，必须先明确共同的源范畴，不能仅因二者的对象取值相同就直接套用以上方块。

\[
\begin{tikzcd}[column sep=large,row sep=large]
F(Y)\arrow[r,"\eta_Y"]\arrow[d,"F(f)"']&
G(Y)\arrow[d,"G(f)"]\\
F(X)\arrow[r,"\eta_X"']&G(X).
\end{tikzcd}
\]

\begin{example}\label{b4:ex:torsion-natural-inclusion}
对交换群 \(A\)，记 \(t(A)=\{a:\;\text{存在整数}\;n\ge1,\ na=0\}\)。
若 \(na=0\)、\(mb=0\)，则 \(nm(a-b)=0\)，所以这是子群。群同态 \(f:A\rightarrow B\) 把挠元送到挠元，因为 \(nf(a)=f(na)=0\)。于是限制映射定义函子 \(t:\mathbf{Ab}\rightarrow\mathbf{Ab}\)，包含
\(j_A:t(A)\hookrightarrow A\) 组成自然变换 \(j:t\Rightarrow\operatorname{Id}\)。
\end{example}

这项自然性没有说挠子群都相同，也没有选取任何补群。它说每个同态都保留这个已经由对象自身确定的子群。同样，对固定整数 \(n\)，逐对象取乘 \(n\) 的同态，得到
\(\operatorname{Id}_{\mathbf{Ab}}\Rightarrow\operatorname{Id}_{\mathbf{Ab}}\)；这里自然性就是 \(f(na)=nf(a)\)。

\subsection{函子复合与自然变换复合}
\label{b4:subsec:0102:02}

函子 \(F:\mathcal C\rightarrow\mathcal D\)、\(H:\mathcal D\rightarrow\mathcal E\) 可以复合，规定
\((HF)(X)=H(F(X))\)、\((HF)(f)=H(F(f))\)。两次使用函子公理，便有
\[
HF(gf)=H\bigl(F(g)F(f)\bigr)=HF(g)\,HF(f).
\]
恒等也被保留。对象和态射都不改变的函子记为 \(\operatorname{Id}_{\mathcal C}\)，复合结合律及恒等律逐对象、逐态射成立。

\begin{proposition}\label{b4:prop:vertical-composition}
设 \(\mathcal C,\mathcal D\) 为范畴，\(F,G,H:\mathcal C\rightarrow\mathcal D\) 为函子，
\(\eta:F\Rightarrow G\)、\(\theta:G\Rightarrow H\) 为自然变换。则
\[
(\theta\eta)_X=\theta_X\eta_X
\]
定义自然变换 \(\theta\eta:F\Rightarrow H\)，称为 \(\eta\) 与 \(\theta\) 的\term[chuizhifuhe]{垂直复合}[vertical composition]，即逐分量复合。这种复合具有结合律；分量为 \(1_{F(X)}\) 的变换 \(1_F\) 为恒等变换。
\end{proposition}
\begin{proof}
对 \(f:X\rightarrow Y\)，依次使用两个自然性等式得到
\[
H(f)\theta_X\eta_X
=\theta_YG(f)\eta_X
=\theta_Y\eta_YF(f).
\]
所以复合自然。三个变换的复合在对象 \(X\) 处是三个 \(\mathcal D\) 态射的复合，结合律及恒等律直接来自 \(\mathcal D\)。
\end{proof}

\begin{definition}\label{b4:def:natural-isomorphism}
设 \(\mathcal C,\mathcal D\) 为范畴，\(F,G:\mathcal C\rightarrow\mathcal D\) 为函子，\(\eta:F\Rightarrow G\) 为自然变换。若存在自然变换 \(\theta:G\Rightarrow F\)，使
\(\theta\eta=1_F\)、\(\eta\theta=1_G\)，称 \(\eta\) 为\term[zirantonggou]{自然同构}[natural isomorphism]，记为 \(F\cong G\)。
\end{definition}

\begin{proposition}\label{b4:prop:natural-iso-components}
设 \(\mathcal C,\mathcal D\) 为范畴，\(F,G:\mathcal C\rightarrow\mathcal D\) 为函子，\(\eta:F\Rightarrow G\) 为自然变换。则 \(\eta\) 为自然同构，当且仅当每个分量 \(\eta_X\)（\(X\in\mathcal C\)）都是 \(\mathcal D\) 中的同构。
\end{proposition}
\begin{proof}
自然逆的各分量显然是分量逆。反过来，假设每个分量可逆。由
\(G(f)\eta_X=\eta_YF(f)\)，左右分别复合相应逆态射，得到
\[
F(f)\eta_X^{-1}=\eta_Y^{-1}G(f).
\]
因此这些逆态射组成一个自然变换 \(G\Rightarrow F\)；逐对象复合为恒等，故是自然逆。逆态射唯一，所以这里不需要为每个对象另作一次选择。
\end{proof}

\begin{example}\label{b4:exm:objectwise-iso-not-natural}
设 \(k\) 为域，\(J\) 为只有对象 \(0,1\) 及一条非恒等态射 \(a:0\rightarrow1\) 的范畴。定义函子 \(F,G:J\rightarrow\mathbf{Vect}_k\)，令
\[
F(0)=F(1)=G(0)=G(1)=k,\qquad
F(a)=1_k,\qquad G(a)=0,
\]
并令两函子把恒等态射送到 \(1_k\)。它们在每个对象处的取值都相同，但若 \(\eta:F\Rightarrow G\) 为自然变换，对 \(a\) 的自然性便给出
\[
0\eta_0=\eta_1 1_k,
\]
所以 \(\eta_1=0\)，不可能是同构。因此 \(F\) 与 \(G\) 不自然同构。函子在对象处的取值既不能决定其态射作用，逐对象存在同构也不能保证存在自然同构。
\end{example}

自然变换还可以与函子在左右两侧复合。若 \(\eta:F\Rightarrow G\)，
\(H:\mathcal D\rightarrow\mathcal E\)，则
\((H\eta)_X=H(\eta_X)\) 给出 \(HF\Rightarrow HG\)，因为
\[
HG(f)\,H(\eta_X)
=H\bigl(G(f)\eta_X\bigr)
=H\bigl(\eta_YF(f)\bigr).
\]
若 \(K:\mathcal B\rightarrow\mathcal C\)，则
\((\eta K)_B=\eta_{K(B)}\) 给出 \(FK\Rightarrow GK\)，自然性就是 \(\eta\) 对箭头 \(K(u)\) 的自然性。

\begin{proposition}\label{b4:prop:horizontal-composition}
设 \(\mathcal C,\mathcal D,\mathcal E\) 为范畴，\(F,G:\mathcal C\rightarrow\mathcal D\)、\(H,L:\mathcal D\rightarrow\mathcal E\) 为函子，
\(\eta:F\Rightarrow G\)、\(\alpha:H\Rightarrow L\) 为自然变换。从 \(HF\) 到 \(LG\)，可以先改变内侧函子，再改变外侧函子，也可以按相反次序复合：
\[
\begin{aligned}
HF&\xRightarrow{H\eta}HG\xRightarrow{\alpha G}LG,\\
HF&\xRightarrow{\alpha F}LF\xRightarrow{L\eta}LG.
\end{aligned}
\]
这两条路径给出同一个自然变换，其分量为
\[
(\alpha*\eta)_X
=\alpha_{G(X)}H(\eta_X)
=L(\eta_X)\alpha_{F(X)}
\]
称这个自然变换为 \(\eta\) 与 \(\alpha\) 的\term[shuipingfuhe]{水平复合}[horizontal composition]，也称横向复合。
\end{proposition}
\begin{proof}
两个公式相等，正是 \(\alpha\) 对 \(\mathcal D\) 中态射
\(\eta_X:F(X)\rightarrow G(X)\) 的自然性。第一个公式表示自然变换
\(H\eta:HF\Rightarrow HG\) 与 \(\alpha G:HG\Rightarrow LG\) 的逐分量复合，故由\cref{b4:prop:vertical-composition}自然。
\[
\begin{tikzcd}[column sep=large,row sep=large]
H(F(X))\arrow[r,"H(\eta_X)"]\arrow[d,"\alpha_{F(X)}"']&
H(G(X))\arrow[d,"\alpha_{G(X)}"]\\
L(F(X))\arrow[r,"L(\eta_X)"']&L(G(X)).
\end{tikzcd}
\]
\end{proof}

这项公式用于比较复合函子。把范畴看作对象、函子看作1态射、自然变换看作2态射，两种复合共同形成严格2范畴的例子。它们之间的交换律及完整结构见\secref{b4:sec:F01}。

\subsection{忘却结构与附加结构}
\label{b4:subsec:0102:03}

忘却函子保留底层映射，只忘掉一部分对象结构。例如
\[
R\text{-}\mathbf{Mod}\longrightarrow\mathbf{Ab}
\longrightarrow\mathbf{Set}
\]
先忘掉 \(R\) 的作用，再忘掉加法。它们确实是函子，因为恒等与复合的底函数都没有改变。反方向则不能把一个任意集合直接宣称为模；必须构造新的元素及运算。

\begin{example}\label{b4:ex:free-module-functor}
固定含幺环 \(R\)。第二卷第3.1节构造的自由左模
\[
R^{(X)}=\left\{\sum_{x\in X}r_xe_x:
r_x\in R,\ \text{只有有限多个}\;r_x\ne0\right\}
\]
给出函子 \(F_R:\mathbf{Set}\rightarrow R\text{-}\mathbf{Mod}\)。
对函数 \(u:X\rightarrow Y\)，规定
\[
F_R(u)\left(\sum_xr_xe_x\right)=\sum_xr_xe_{u(x)}.
\]
若不同 \(x\) 有同一个像，就把相应系数相加；支集有限保证和有意义。这个公式保持加法与左标量作用。又有
\(F_R(1_X)(e_x)=e_x\) 及
\(F_R(vu)(e_x)=e_{v(u(x))}=F_R(v)F_R(u)(e_x)\)，由基生成性可知它保持恒等与复合。
\end{example}
\symidx{\(R^{(X)}\)：以集合为基的自由左模}
\symidx{\(F_R\)：自由模函子}

设 \(U:R\text{-}\mathbf{Mod}\rightarrow\mathbf{Set}\) 为忘却函子。映射
\(\iota_X:X\rightarrow U F_R(X)\)、\(x\mapsto e_x\)，组成自然变换
\(\operatorname{Id}_{\mathbf{Set}}\Rightarrow UF_R\)，因为
\(UF_R(u)(e_x)=e_{u(x)}\)。另一个方向有线性映射
\[
\epsilon_M:F_RU(M)\longrightarrow M,\qquad
\sum_{m\in M}r_me_m\longmapsto\sum_mr_mm.
\]
对模同态 \(f:M\rightarrow N\)，在基元 \(e_m\) 上，
\(f\epsilon_M(e_m)=f(m)=\epsilon_NF_RU(f)(e_m)\)，所以这些映射也自然。这两族映射并非一对自然逆。虽然
\(U(\epsilon_M)\iota_{U(M)}=1_{U(M)}\)，但当 \(R\ne0\) 时，\(F_RU(R)\) 中对应零元素的基元 \(e_0\) 是非零向量，而 \(\epsilon_R(e_0)=0\)。因此 \(\epsilon_R\) 不单射，不能是同构；\(e_0\) 与自由模中的零向量必须区别。第3章会说明这两族映射怎样组成自由与忘却之间的伴随。

\ProseParagraphBreak
构造是否能够成为函子，还取决于允许哪些态射。例如群的中心 \(Z(G)\) 对同构有相容的映射，但任意同态
\(f:G\rightarrow H\) 未必把 \(Z(G)\) 送入 \(Z(H)\)。取
\(G=\{1,(12)\}\subset S_3=H\)，则 \(G\) 交换，中心为自身，而 \(S_3\) 的中心平凡；包含映射便不能限制为 \(Z(G)\rightarrow Z(H)\)。问题不是中心不存在，而是所提议的态射规则不成立。

% !TeX root = ../../Book4.tex
% 纲要标题：### 1.3 范畴等价与常见范畴构造
% 纲要位置：计划纲要.md，第 2463 行。

\section{范畴等价与常见范畴构造}
\label{b4:sec:0103}

判断一个函子是否保留了原有结构，不能只数对象的同构类型，还要比较每对对象之间的态射。对偶范畴、函子范畴和切片范畴提供了改变比较方式的具体例子。

% 纲要内容：
%
% * 全忠实性、本质满性与范畴等价
% * 对偶范畴与积范畴
% * 函子范畴与箭头范畴
% * 切片、余切片与逗号范畴
%

\subsection{全忠实性、本质满性与范畴等价}
\label{b4:subsec:0103:01}

一个函子可能忘掉对象的一部分结构，也可能只为同一结构换一种表示。要区分这两种情况，既要检查对象是否都被包括，又要检查对象之间的每个态射是否得到保留。只比较对象的同构类，会漏掉后一个问题。

\begin{definition}\label{b4:def:fully-faithful-essentially-surjective}
设 \(F:\mathcal C\rightarrow\mathcal D\) 为局部小范畴之间的函子。对每对对象 \(X,Y\)，考虑函数
\[
F_{X,Y}:\operatorname{Hom}_{\mathcal C}(X,Y)\longrightarrow\operatorname{Hom}_{\mathcal D}(FX,FY),
\qquad f\longmapsto F(f).
\]
若这些函数都满射，称 \(F\) 为\term[quanhanzi]{全函子}[full functor]；若都单射，称为\term[zhongshihanzi]{忠实函子}[faithful functor]；若都双射，称为\term[quanzhongshihanzi]{全忠实函子}[fully faithful functor]。若每个 \(Y\in\mathcal D\) 都同构于某个 \(F(X)\)，称 \(F\) 为\term[benzhimanhanzi]{本质满函子}[essentially surjective functor]。
\end{definition}

\begin{definition}\label{b4:def:subcategory}
设 \(\mathcal C\) 为范畴。选择其中一部分对象，并在这些对象之间选择包含所有相应恒等态射、且对复合封闭的一部分态射，采用原来的复合，得到的范畴称为 \(\mathcal C\) 的子范畴。若所选任意对象 \(X,Y\) 之间的全部 \(\operatorname{Hom}_{\mathcal C}(X,Y)\) 都被保留，称为\term[quanzifanchou]{全子范畴}[full subcategory]。
\end{definition}

子范畴的包含函子忠实，全子范畴的包含函子全忠实。规定一类对象后，还须说明是否保留它们之间的全部态射；例如有限维空间的全子范畴，与只允许这些空间之间同构的子范畴，是两个不同范畴。

\ProseParagraphBreak
“忠实”只比较固定起点和终点的两条箭头，不要求对象映射单射。例如，源范畴有两个对象，每对对象间恰有一个箭头；目标范畴只有一个对象及其恒等箭头。把两个对象都送到这个对象、把所有箭头都送到恒等箭头，便得到函子。每个 Hom 集上的映射都是单点集之间的双射，所以它全忠实，尽管合并了对象；\cref{b4:ex:equivalent-not-isomorphic}还将写出它的拟逆。反过来，给两个不同对象取不同的像，并不能保证它们之间的态射没有被合并。

\ProseParagraphBreak
忘却函子 \(\mathbf{Vect}_k\rightarrow\mathbf{Set}\) 忠实，但不全，因为函数 \(k\rightarrow k\)、\(x\mapsto x+1\) 不是线性映射。包含
\(\mathbf{Ab}\rightarrow\mathbf{Grp}\) 全忠实，但不本质满，因为非交换群不能同构于交换群。对偏序集之间的保序映射，所得函子总忠实；它全当且仅当
\(F(x)\le F(y)\) 还能推出 \(x\le y\)。这些条件各自检测不同的信息。

\begin{proposition}\label{b4:prop:fully-faithful-reflects-isomorphisms}
设 \(\mathcal C,\mathcal D\) 为局部小范畴，\(F:\mathcal C\rightarrow\mathcal D\) 为全忠实函子。若 \(\mathcal C\) 中态射 \(f:X\rightarrow Y\) 的像 \(F(f)\) 可逆，则 \(f\) 可逆；对 \(\mathcal C\) 中任意对象 \(X,Y\)，若 \(FX,FY\) 同构，则 \(X,Y\) 同构。
\end{proposition}
\begin{proof}
全性给出 \(g:Y\rightarrow X\)，使 \(F(g)=F(f)^{-1}\)。因此
\(F(gf)=1_{FX}=F(1_X)\)，由忠实性得 \(gf=1_X\)；另一侧同理。若只给 \(a:FX\rightarrow FY\) 为同构，先用全性提升为 \(f:X\rightarrow Y\)，再用已证结论。
\end{proof}

\begin{definition}\label{b4:def:equivalence}
设 \(\mathcal C,\mathcal D\) 为范畴，\(F:\mathcal C\rightarrow\mathcal D\) 为函子。若存在函子
\(G:\mathcal D\rightarrow\mathcal C\) 及自然同构
\[
\eta:\operatorname{Id}_{\mathcal C}\Rightarrow GF,\qquad
\varepsilon:FG\Rightarrow\operatorname{Id}_{\mathcal D},
\]
称 \(F\) 为\term[fanchoudengjia]{范畴等价}[equivalence of categories]，称 \(G\) 为其\term[nini]{拟逆}[quasi-inverse]，记 \(\mathcal C\simeq\mathcal D\)。若存在 \(G\) 使
\(GF=\operatorname{Id}_{\mathcal C}\)、\(FG=\operatorname{Id}_{\mathcal D}\) 作为函子严格相等，称 \(F\) 为范畴同构。
\end{definition}
\symidx{\(\operatorname{Id}_{\mathcal C}\)：范畴的恒等函子}
\symidx{\(\mathcal C\simeq\mathcal D\)：范畴等价}

范畴同构在对象上必须给出真正的双射；范畴等价允许把互相同构的不同对象换成一个代表。自然同构保证这种替换对全部态射相容，而不只是在每个对象处分别碰巧有一个同构。

\begin{theorem}[范畴等价的判据]\label{b4:thm:equivalence-criterion}
设 \(\mathcal C,\mathcal D\) 为局部小范畴，\(F:\mathcal C\rightarrow\mathcal D\) 为函子。
\begin{enumerate}
\item 若 \(F\) 是等价，则它全忠实且本质满。
\item 若 \(F\) 全忠实，并且已为每个 \(Y\in\mathcal D\) 指定对象 \(X_Y\in\mathcal C\) 及同构 \(\varepsilon_Y:F(X_Y)\rightarrow Y\)，则这些数据给出 \(F\) 的一个拟逆。
\end{enumerate}
特别地，对于小范畴，在选择公理下，\(F\) 是等价当且仅当它全忠实且本质满。
\end{theorem}
\begin{proof}
先设 \(G,\eta,\varepsilon\) 为等价数据。每个
\(\varepsilon_Y:F(GY)\rightarrow Y\) 都是同构，所以取 \(X_Y=GY\) 即见 \(F\) 本质满。若
\(F(f)=F(g)\)，其中 \(f,g:X\rightarrow X'\)，自然性给出
\[
\eta_{X'}f=GF(f)\eta_X=GF(g)\eta_X=\eta_{X'}g.
\]
消去同构 \(\eta_{X'}\)，得 \(f=g\)，所以 \(F\) 忠实。若
\(u,v:Y\rightarrow Y'\) 满足 \(G(u)=G(v)\)，则由 \(\varepsilon\) 的自然性，
\[
u\varepsilon_Y
=\varepsilon_{Y'}FG(u)
=\varepsilon_{Y'}FG(v)
=v\varepsilon_Y.
\]
消去同构 \(\varepsilon_Y\)，得 \(u=v\)，故 \(G\) 也忠实。

为证明 \(F\) 的全性，给定 \(a:FX\rightarrow FX'\)。先用 \(G\) 得到
\(G(a):GFX\rightarrow GFX'\)，再用 \(\eta_X\) 和 \(\eta_{X'}^{-1}\) 把端点改回 \(X,X'\)，便得到
\(f=\eta_{X'}^{-1}G(a)\eta_X\)。对这个 \(f\) 使用 \(\eta\) 的自然性，得
\[
GF(f)\eta_X=\eta_{X'}f=G(a)\eta_X.
\]
约去 \(\eta_X\)，再用 \(G\) 的忠实性，得到 \(F(f)=a\)。因此 \(F\) 全。

反向，令 \(G(Y)=X_Y\)。对 \(a:Y\rightarrow Y'\)，全忠实性保证存在唯一态射 \(G(a)\)，使
\begin{equation}\label{b4:eq:quasi-inverse-arrow}
F(G(a))=\varepsilon_{Y'}^{-1}a\varepsilon_Y.
\end{equation}
代入恒等箭头，右侧为恒等，所以 \(G(1_Y)=1_{G(Y)}\)。对可复合的 \(a,b\)，右侧共轭公式满足
\[
\bigl(\varepsilon_{Y''}^{-1}b\varepsilon_{Y'}\bigr)
\bigl(\varepsilon_{Y'}^{-1}a\varepsilon_Y\bigr)
=\varepsilon_{Y''}^{-1}ba\varepsilon_Y.
\]
由 \(F\) 忠实，\(G(b)G(a)=G(ba)\)。故 \(G\) 是函子，且\eqref{b4:eq:quasi-inverse-arrow}正是 \(\varepsilon:FG\Rightarrow\operatorname{Id}_{\mathcal D}\) 的自然性。

对 \(X\in\mathcal C\)，用全忠实性唯一确定
\(\eta_X:X\rightarrow GFX\)，满足
\[
F(\eta_X)=\varepsilon_{FX}^{-1}.
\]
由\cref{b4:prop:fully-faithful-reflects-isomorphisms}，\(\eta_X\) 可逆。若 \(f:X\rightarrow X'\)，\(\varepsilon\) 对 \(F(f)\) 的自然性给出
\[
F(f)\varepsilon_{FX}
=\varepsilon_{FX'}FGF(f).
\]
复合相应逆态射，得到
\[
\begin{aligned}
F\bigl(GF(f)\eta_X\bigr)
&=FGF(f)\varepsilon_{FX}^{-1}\\
&=\varepsilon_{FX'}^{-1}F(f)
=F\bigl(\eta_{X'}f\bigr).
\end{aligned}
\]
忠实性给出 \(\eta\) 的自然性，故得到所需等价。

当两个范畴都小时，本质满使每个 \(Y\) 的候选对
\((X,\varepsilon:FX\cong Y)\) 构成非空集合。这些集合由
\(\operatorname{Ob}\mathcal D\) 这个集合取指标，选择公理给出第二项所需的数据。
\end{proof}

对大范畴，判据第二项把所需选择明确放在假设中；单独说“本质满”只给出逐对象的存在命题，不会自动提供真类指标的选择函数。Weibel 关于从每个同构类选取代表对象及相应同构的讨论也明确以这些选择为前提，见\cite[A.2.4, p. 422; A.3.2, pp. 423--424]{Weibel1994HomologicalAlgebra}。

\begin{example}\label{b4:ex:equivalent-not-isomorphic}
令 \(\mathcal E\) 有两个对象 \(a,b\)，每对对象之间恰有一个态射。所有复合由唯一性决定，箭头均可逆。令 \(\mathbf1\) 为只有一个对象及恒等态射的范畴。唯一函子
\(\mathcal E\rightarrow\mathbf1\) 全忠实且本质满，反向选取 \(a\) 就给出一个具体拟逆。它们等价，却不可能同构，因为对象数不同。
\end{example}

\subsection{对偶范畴与积范畴}
\label{b4:subsec:opposite-product-categories}
\label{b4:subsec:0103:02}

箭头反向以后，同一套范畴公理把单态射与满态射等概念配成对偶；把两个范畴的对象和态射逐对放在一起，则能同时记录两个变量。对偶范畴和积范畴结合起来，尤其适合描述 Hom 这样的双变量构造。

\begin{definition}\label{b4:def:opposite-category}
设 \(\mathcal C\) 为范畴。保留同一对象类，规定
\(\operatorname{Hom}_{\mathcal C^{\mathrm{op}}}(X,Y)=\operatorname{Hom}_{\mathcal C}(Y,X)\)，并把复合反向：若
\(f^{\mathrm{op}}:Y\rightarrow X\)、\(g^{\mathrm{op}}:Z\rightarrow Y\) 分别来自
\(f:X\rightarrow Y\)、\(g:Y\rightarrow Z\)，则
\[
f^{\mathrm{op}}\circ g^{\mathrm{op}}=(g\circ f)^{\mathrm{op}}.
\]
所得范畴称为 \(\mathcal C\) 的\term[duioufanchou]{对偶范畴}[opposite category]。
\end{definition}
\symidx{\(\mathcal C^{\mathrm{op}}\)：相反范畴}

原范畴的结合律与恒等律反向后仍成立，所以这是一个范畴；再取对偶恢复原范畴。一个从 \(\mathcal C\) 到 \(\mathcal D\) 的反变函子，正是一个协变函子
\(\mathcal C^{\mathrm{op}}\rightarrow\mathcal D\)。单态射与满态射的定义也互为对偶，因为左消去与右消去在箭头反向后交换。

\begin{definition}\label{b4:def:product-category}
设 \(\mathcal C,\mathcal D\) 为范畴。以对象对 \((X,Y)\) 为对象，规定
\[
\operatorname{Hom}_{\mathcal C\times\mathcal D}\bigl((X,Y),(X',Y')\bigr)
=\operatorname{Hom}_{\mathcal C}(X,X')\times\operatorname{Hom}_{\mathcal D}(Y,Y'),
\]
恒等与复合逐分量计算，所得范畴称为它们的积范畴。
\end{definition}
\symidx{\(\mathcal C\times\mathcal D\)：积范畴}

两个投影是函子。积范畴使一个同时依赖多个对象的构造仍可写成普通函子。例如，对局部小范畴 \(\mathcal C\)，有
\[
\operatorname{Hom}_{\mathcal C}:
\mathcal C^{\mathrm{op}}\times\mathcal C\longrightarrow\mathbf{Set}.
\]
箭头 \((a^{\mathrm{op}},b):(X,Y)\rightarrow(X',Y')\)，其中
\(a:X'\rightarrow X\)、\(b:Y\rightarrow Y'\)，把 \(u:X\rightarrow Y\) 送到 \(bua\)。复合规律来自结合律；第一个变量反变、第二个变量协变已经由源范畴中的对偶符号体现。

\subsection{函子范畴与箭头范畴}
\label{b4:subsec:functor-arrow-categories}

自然变换可以复合，也有恒等变换，因此函子本身可以成为一个新范畴的对象。若源范畴只记录一条箭头，这个构造就把原范畴中的箭头作为对象，把交换方块作为它们之间的态射。

\begin{definition}\label{b4:def:functor-category}
设 \(\mathcal C\) 为小范畴，\(\mathcal D\) 为局部小范畴。以函子
\(F:\mathcal C\rightarrow\mathcal D\) 为对象、自然变换为态射，并按分量复合，所得范畴记为
\(\operatorname{Fun}(\mathcal C,\mathcal D)\)，也可以记为 \([\mathcal C,\mathcal D]\) 或 \(\mathcal D^{\mathcal C}\)，称为\term[hanzifanchou]{函子范畴}[functor category]。对其中的函子 \(F,G\)，其态射集合记为 \(\operatorname{Nat}(F,G)\)，即
\[
\operatorname{Hom}_{\operatorname{Fun}(\mathcal C,\mathcal D)}(F,G)
=\operatorname{Nat}(F,G).
\]
\end{definition}
\symidx{\(\operatorname{Fun}(\mathcal C,\mathcal D),\ [\mathcal C,\mathcal D],\ \mathcal D^{\mathcal C}\)：函子范畴}
\symidx{\(\operatorname{Nat}(F,G)\)：自然变换集合}

\begin{proposition}\label{b4:prop:functor-category-locally-small}
设 \(\mathcal C\) 为小范畴、\(\mathcal D\) 为局部小范畴，则
\(\operatorname{Fun}(\mathcal C,\mathcal D)\) 是局部小范畴；若 \(\mathcal D\) 也小，它还是小范畴。
\end{proposition}
\begin{proof}
复合、恒等与范畴公理由\cref{b4:prop:vertical-composition}得到。固定 \(F,G\) 后，所有分量族属于集合
\[
\prod_{X\in\operatorname{Ob}\mathcal C}\operatorname{Hom}_{\mathcal D}(FX,GX),
\]
而自然变换正是其中满足所有自然性等式的子集，因此 Hom 为集合。若 \(\mathcal D\) 也小，对象映射与态射映射分别是两个集合之间的函数；所有满足函子条件的这样的函数对仍组成集合，故函子范畴小。
\end{proof}

源范畴小保证上式是集合指标的乘积。若源对象组成真类，这个乘积未必是集合，因而不能据此推出函子范畴局部小。

\ProseParagraphBreak
令 \([1]\) 为两个对象 \(0,1\) 及一个非恒等箭头 \(0\rightarrow1\) 组成的范畴。函子
\([1]\rightarrow\mathcal C\) 恰好指定 \(\mathcal C\) 中的一条箭头；两者之间的自然变换恰好是一个交换方块。因此
\(\operatorname{Fun}([1],\mathcal C)\) 是箭头范畴，它把“箭头之间的映射”变成一个可以继续研究的对象。

\subsection{切片、余切片与逗号范畴}
\label{b4:subsec:slice-coslice-comma-categories}
\label{b4:subsec:0103:03}

箭头范畴允许一条箭头的两个端点都变化。如果固定其中一端，并要求所有比较都保持这个端点，就得到切片或余切片范畴。指定映射不仅附着在原对象上，还决定了新范畴中哪些态射可以使用。

\begin{definition}\label{b4:def:slice-coslice-category}
设 \(\mathcal C\) 为范畴，\(A\in\mathcal C\)。
以态射 \(p:X\rightarrow A\) 为对象，以满足 \(qh=p\) 的态射
\(h:X\rightarrow Y\) 作为 \(p:X\rightarrow A\) 到 \(q:Y\rightarrow A\) 的态射，得到\term[qiepianfanchou]{切片范畴}[slice category] \(\mathcal C/A\)。
若改以态射 \(p:A\rightarrow X\) 为对象，以满足 \(hp=q\) 的态射 \(h:X\rightarrow Y\) 作为 \(p:A\rightarrow X\) 到 \(q:A\rightarrow Y\) 的态射，仍用 \(\mathcal C\) 的恒等与复合，则得到余切片范畴 \(A/\mathcal C\)。
\end{definition}
\symidx{\(\mathcal C/A,\ A/\mathcal C\)：切片范畴与余切片范畴}

\[
\begin{tikzcd}[column sep=large,row sep=large]
X\arrow[r,"h"]\arrow[dr,"p"']&Y\arrow[d,"q"]\\
&A
\end{tikzcd}
\qquad
\begin{tikzcd}[column sep=large,row sep=large]
&A\arrow[dl,"p"']\arrow[d,"q"]\\
X\arrow[r,"h"']&Y.
\end{tikzcd}
\]

恒等态射满足所列三角形条件。若 \(qh=p\)、\(rk=q\)，则
\(r(kh)=(rk)h=qh=p\)，所以切片中的复合仍在其中，结合律来自 \(\mathcal C\)；余切片同理。两者的 Hom 都是原 Hom 集的子集，因而在 \(\mathcal C\) 局部小时局部小。例子是
\(\{*\}/\mathbf{Set}\)：从单点集到 \(X\) 的映射指定一个点，三角形条件要求态射保持这个点，所以它就是带基点集合范畴。

\ProseParagraphBreak
切片固定一个对象，再记录到它或从它出发的映射。有时映射两端都来自构造：例如给一个集合 \(X\)，研究把它映到某个模的底集合中的所有方式。逗号范畴把这些“对象连同一个指定映射”统一写出。

\begin{definition}\label{b4:def:comma-category}
设 \(\mathcal A,\mathcal B,\mathcal C\) 为范畴，
\(S:\mathcal A\rightarrow\mathcal C\)、\(T:\mathcal B\rightarrow\mathcal C\) 为函子。
以三元组
\((A,B,u)\)，其中 \(A\in\mathcal A\)、\(B\in\mathcal B\)、\(u:S(A)\rightarrow T(B)\)，为对象。态射
\[
(a,b):(A,B,u)\longrightarrow(A',B',u')
\]
由 \(a:A\rightarrow A'\)、\(b:B\rightarrow B'\) 组成，要求
\[
T(b)u=u'S(a).
\]
按两个分量分别复合，所得范畴记为 \((S\downarrow T)\)，称为\term[douhaofanchou]{逗号范畴}[comma category]。
\end{definition}
\symidx{\((S\downarrow T)\)：逗号范畴}

\[
\begin{tikzcd}[column sep=large,row sep=large]
S(A)\arrow[r,"u"]\arrow[d,"S(a)"']&T(B)\arrow[d,"T(b)"]\\
S(A')\arrow[r,"u'"']&T(B').
\end{tikzcd}
\]

\begin{proposition}\label{b4:prop:comma-category-construction}
设 \(\mathcal A,\mathcal B,\mathcal C\) 为范畴，\(S:\mathcal A\rightarrow\mathcal C\)、\(T:\mathcal B\rightarrow\mathcal C\) 为函子。\cref{b4:def:comma-category}给出的逗号构造 \((S\downarrow T)\) 在逐分量的恒等与复合下构成范畴，向 \(\mathcal A,\mathcal B\) 的两个投影是函子。若
\(\mathcal A,\mathcal B\) 局部小，则 \((S\downarrow T)\) 局部小。
\end{proposition}
\begin{proof}
恒等对满足方块条件。若先后有 \((a,b)\)、\((a',b')\)，则
\[
T(b'b)u=T(b')T(b)u
=T(b')u'S(a)
=u''S(a')S(a)=u''S(a'a).
\]
所以复合仍满足条件。结合律及恒等律逐分量继承；投影显然保持它们。固定两个三元组后，态射集合是
\(\operatorname{Hom}_{\mathcal A}(A,A')\times\operatorname{Hom}_{\mathcal B}(B,B')\) 的一个子集，故局部小。
\end{proof}

令 \(\mathbf1\rightarrow\mathcal C\) 选择对象 \(A\)。取 \(S=\operatorname{Id}_{\mathcal C}\)，便得到
\((\operatorname{Id}_{\mathcal C}\downarrow A)=\mathcal C/A\)；取
\(T=\operatorname{Id}_{\mathcal C}\)，则得到
\((A\downarrow\operatorname{Id}_{\mathcal C})=A/\mathcal C\)。这里省略了三元组中固定不变的分量。

\ProseParagraphBreak
对忘却函子 \(U:R\text{-}\mathbf{Mod}\rightarrow\mathbf{Set}\) 和固定集合 \(X\)，仍把选择 \(X\) 的函子 \(\mathbf1\rightarrow\mathbf{Set}\) 简记为 \(X\)。于是
\((X\downarrow U)\) 的对象是一个模 \(M\) 及函数 \(j:X\rightarrow U(M)\)。
从 \((M,j)\) 到 \((N,k)\) 的态射是满足 \(U(f)j=k\) 的线性映射 \(f:M\rightarrow N\)。
自由模给出对象 \((F_R(X),\iota_X)\)，且第二卷第3.1节的泛性质说明，从它到每个 \((M,j)\) 恰有一个态射。下一章把这种性质命名为始对象；在这里，指定函数 \(j\) 已经被纳入对象，所以唯一性的条件不再需要附加在范畴之外。

\ProseParagraphBreak
另一个例子固定两个对象 \(X,Y\in\mathcal C\)。令离散的单对象范畴分别选出它们，则逗号范畴的对象是态射 \(X\rightarrow Y\)，而只有恒等态射。相比之下，若取
\((\operatorname{Id}_{\mathcal C}\downarrow\operatorname{Id}_{\mathcal C})\)，则两端都可变化，得到的是箭头范畴。是否固定端点，决定了它们之间允许哪些比较。

% !TeX root = ../../Book4.tex
% 纲要标题：### 1.4 自然同构与典范同构
% 纲要位置：计划纲要.md，第 2469 行。

\section{自然同构与典范同构}
\label{b4:sec:0104}

同构的存在与同构的自然性是两件事。张量积、对偶及模分解中的映射，能说明哪些对应由结构决定，哪些对应仍取决于选基或选补模。

% 纲要内容：
%
% * 双对偶与张量积的自然同构
% * 选基同构与典范同构
% * 挠子模、无挠商与有限生成模的非自然分裂
%
% ---
%

\subsection{双对偶与张量积的自然同构}
\label{b4:subsec:0104:01}

第二卷已经在模的具体构造中使用自然性。本节用函子和自然变换描述这些映射，辨认其源范畴、目标范畴和分量。

\ProseParagraphBreak
自然性最初也来自一个具体困难：一族对象逐个同构，是否就能在取极限后得到同构？Eilenberg 与 Mac Lane 的一般理论起于研究同调中的这种问题。\footnote{取极限的研究动机见\cite[Introduction, p. 236, especially footnote 4]{EilenbergMacLane1945NaturalEquivalences}；该脚注说明，从复形同调转向空间同调时需要这种极限过程。}逐项同构还须与系统中的过渡映射交换，才是两套系统之间的同构；这种相容性正是自然性。他们以有限维实向量空间为例，对比选基得到的 \(V\cong V^*\) 与求值映射给出的 \(V\cong V^{**}\)：前者随基改变，后者同时与所有线性映射相容。\footnote{对偶与双对偶的例子见\cite[Introduction, pp. 232--234]{EilenbergMacLane1945NaturalEquivalences}。}下面把求值映射的自然性与有限维时的可逆性分别写清楚。

\ProseParagraphBreak
设 \(k\) 为域。两次对偶给出协变函子
\(D^2:\mathbf{Vect}_k\rightarrow\mathbf{Vect}_k\)，其中
\(D^2(V)=V^{**}\)、\(D^2(f)=f^{**}\)。求值映射
\[
j_V:V\longrightarrow V^{**},\qquad j_V(v)(\lambda)=\lambda(v)
\]
组成 \(\operatorname{Id}\Rightarrow D^2\) 的自然变换。第二卷定理1.2.9“有限维双对偶同构”已经证明
\[
f^{**}j_V=j_Wf
\]
对任意线性映射 \(f:V\rightarrow W\) 成立，并证明有限维时 \(j_V\) 可逆。结合\cref{b4:prop:natural-iso-components}，它在
\(\mathbf{Vect}_k^{\mathrm{fd}}\) 上是自然同构。有限维保证分量的满射性，自然性本身并不需要这个条件。

\[
\begin{tikzcd}[column sep=large,row sep=large]
V\arrow[r,"j_V"]\arrow[d,"f"']&V^{**}\arrow[d,"f^{**}"]\\
W\arrow[r,"j_W"']&W^{**}.
\end{tikzcd}
\]

再设 \(R\) 为交换含幺环，\(L,M,N\) 为幺性 \(R\) 模。第二卷定理6.2.3“结合与单位同构”已经构造
\[
\begin{aligned}
\lambda_M &:R\otimes_RM\longrightarrow M,
&r\otimes m&\longmapsto rm,\\
\alpha_{L,M,N}&:(L\otimes_RM)\otimes_RN
\longrightarrow L\otimes_R(M\otimes_RN),
&(l\otimes m)\otimes n&\longmapsto l\otimes(m\otimes n).
\end{aligned}
\]
第一族是两个从 \(R\text{-}\mathbf{Mod}\) 到自身的函子之间的自然同构，第二族则从三个模范畴的积范畴出发。若分别给三个变量施以同态 \(f,g,h\)，两种复合在纯张量上的值都是
\(f(l)\otimes(g(m)\otimes h(n))\)，因此三变量自然性有明确的可核验含义。此处采用交换环以免另记外侧双模作用；一般双模版本仍以第二卷的陈述为准。

\ProseParagraphBreak
这两个例子也说明一个同构族必须说清楚相对于哪些态射自然。双对偶涉及所有线性映射，张量结合涉及三个变量上的同态。只说“公式没有写基”可以帮助猜测自然性，却不能代替相应方块的等式。

\subsection{选基同构与典范同构}
\label{b4:subsec:0104:02}

通常把由给定数据直接确定、或由明确的相容条件唯一刻画的映射称为典范映射。“典范”需要指出相对于哪份数据；自然性则是给定两个函子以后可以检验的一组等式。两者有密切联系，但不是同一个定义。

\begin{proposition}\label{b4:prop:natural-scalars}
设 \(k\) 为域。对向量空间及全部线性映射组成的范畴 \(\mathbf{Vect}_k\)，每个自然变换
\(\eta:\operatorname{Id}_{\mathbf{Vect}_k}\Rightarrow\operatorname{Id}_{\mathbf{Vect}_k}\)
都唯一具有形式
\[
\eta_V=c\,1_V\qquad(V\in\mathbf{Vect}_k),
\]
其中 \(c\in k\)。它是自然同构当且仅当 \(c\ne0\)。
\end{proposition}
\begin{proof}
令 \(c=\eta_k(1)\)。由于 \(\eta_k\) 线性，它为乘 \(c\) 的映射。任取 \(v\in V\)，线性映射 \(a_v:k\rightarrow V\)、\(r\mapsto rv\)，使自然性给出
\[
\eta_V(v)=\eta_Va_v(1)=a_v\eta_k(1)=cv.
\]
所以 \(c\) 决定全部分量，并由 \(V=k\) 保证唯一。反过来，标量乘法与每个线性映射交换，所以确实自然。各分量可逆当且仅当在 \(k\) 上可逆，也就是 \(c\ne0\)。
\end{proof}

因此即使已经要求自然，同构族也未必唯一：每个非零 \(c\) 都给出一个自然自同构。若另要求在 \(k\) 上把 \(1\) 送到 \(1\)，才得到恒等族。使用“唯一的典范同构”时，必须说出这样的确定条件，或给出相关泛性质。

\begin{example}\label{b4:ex:basis-dual-not-natural}
在实直线 \(V=\mathbb R e\) 中，以基 \(e\) 规定同构
\(b_e:V\rightarrow V^*\)、\(e\mapsto e^*\)。改用基 \(e'=2e\)，则其对偶基为
\((e')^*=\dfrac12e^*\)，所以
\[
b_{e'}(e)=\dfrac14e^*,\qquad b_e(e)=e^*.
\]
选基产生了同构，却没有产生同一个映射。

\ProseParagraphBreak
更严格地说，普通对偶反变，不能把 \(\operatorname{Id}\) 和 \(D\) 直接作为两个同向函子。令 \(\mathcal V_{\mathrm{iso}}\) 为有限维实向量空间及其间线性同构组成的范畴。在这里，对偶可以按下列协变规则重新组织：
\[
D_{\mathrm{iso}}:\mathcal V_{\mathrm{iso}}\longrightarrow\mathcal V_{\mathrm{iso}},
\qquad V\longmapsto V^*,\quad f\longmapsto(f^{-1})^*.
\]
假如存在自然同构
\(b:\operatorname{Id}_{\mathcal V_{\mathrm{iso}}}\Rightarrow D_{\mathrm{iso}}\)，对非零有限维实空间的同构
\(f=2\,1_V\)，自然性将要求
\[
\dfrac12 b_V=b_V(2\,1_V)=2b_V,
\]
故 \(b_V=0\)，矛盾。这个反例针对的是这两个确定的函子及所有实线性同构。
\end{example}

不过，不能从这个例子推出“使用选择的同构都不自然”。如果连同函子对态射的规则也一起改变，选定的同构反而可以组成自然同构。下面写清这一步。

\begin{proposition}\label{b4:prop:transported-functor}
设 \(\mathcal C,\mathcal D\) 为范畴，\(F:\mathcal C\rightarrow\mathcal D\) 为函子。假定对每个 \(X\in\mathcal C\) 已指定对象 \(G_X\in\mathcal D\) 及同构
\(\theta_X:F(X)\rightarrow G_X\)。则唯一存在函子
\(G:\mathcal C\rightarrow\mathcal D\)，其对象值为 \(G_X\)，且这些 \(\theta_X\) 组成自然同构 \(F\Rightarrow G\)。它的态射作用为
\[
G(f)=\theta_YF(f)\theta_X^{-1}
\qquad(f:X\rightarrow Y).
\]
\end{proposition}
\begin{proof}
自然性 \(G(f)\theta_X=\theta_YF(f)\) 已经强制给出所列公式，所以若存在就唯一。按公式定义后，恒等箭头映为恒等；对可复合 \(f,g\)，中间的
\(\theta_Y^{-1}\theta_Y\) 消去，得到 \(G(g)G(f)=G(gf)\)。公式本身又给出自然性，分量可逆保证自然同构。
\end{proof}

例如，设 \(k\) 为域，\(\mathcal V\) 是由一族有限维 \(k\) 向量空间及其间全部线性映射组成的全子范畴。记包含函子为
\(J:\mathcal V\hookrightarrow\mathbf{Vect}_k^{\mathrm{fd}}\)，并为每个 \(V\in\mathcal V\) 指定一组基，得到坐标同构
\(\theta_V:V\rightarrow k^{\dim V}\)。上一命题给出函子
\[
\begin{aligned}
G&:\mathcal V\longrightarrow\mathbf{Vect}_k^{\mathrm{fd}},\\
G(V)&=k^{\dim V},\qquad
G(f)=\theta_Wf\theta_V^{-1}\quad(f:V\longrightarrow W).
\end{aligned}
\]
其态射作用就是相应的矩阵规则，而 \(\theta:J\Rightarrow G\) 是自然同构。标准坐标空间不必属于 \(\mathcal V\)，所以这里比较的是两个取值于 \(\mathbf{Vect}_k^{\mathrm{fd}}\) 的函子。所选基没有消失，而是进入了 \(G\) 的定义；这些基族作为数据给定，其大小约定见\subsecref{b4:subsec:0101:03}。

\subsection{模结构定理中的自然性}
\label{b4:subsec:0104:03}

第二卷定理4.3.2已经证明：主理想整环 \(R\) 上的有限生成模具有不变因子分解
\[
M\cong R^r\oplus R/(d_1)\oplus\cdots\oplus R/(d_s),
\qquad d_1\mid\cdots\mid d_s,
\]
其中 \(d_i\) 为非零非单位。定理4.3.4进一步证明，自由秩和不变因子的相伴类由 \(M\) 的同构类型唯一确定。这并没有为每个 \(M\) 指定一个唯一分解同构；要讨论分解随同态怎样变化，还须分别考察其中的子模、商模与分裂映射。

\begin{proposition}\label{b4:prop:torsion-quotient-functor}
设 \(R\) 为交换整环。对每个幺性左 \(R\) 模 \(M\)，令
\[
t_R(M)=\{m\in M:\text{存在 }0\ne a\in R,\ am=0\},
\qquad Q_R(M)=M/t_R(M).
\]
则 \(t_R,Q_R\) 都是 \(R\text{-}\mathbf{Mod}\) 上的函子。包含
\(t_R(M)\hookrightarrow M\) 与商映射 \(M\rightarrow Q_R(M)\) 都随 \(M\) 自然。
\end{proposition}
\begin{proof}
若 \(am=0\)、\(bn=0\)，其中 \(a,b\ne0\)，则
\(ab(m-n)=0\)，且 \(ab\ne0\)；标量倍也仍被 \(a\) 零化，所以 \(t_R(M)\) 是子模。对同态 \(f:M\rightarrow N\)，有
\(a\cdot f(m)=f(am)=0\)，故 \(f(t_R(M))\subseteq t_R(N)\)。限制与诱导商映射于是良定义；它们对恒等与复合的相容性直接由限制和商的定义得到。包含方块显然交换；商方块在元素 \(m+t_R(M)\) 上的两个像都为 \(f(m)+t_R(N)\)。
\end{proof}

当 \(R\) 为 PID、\(M\) 有限生成时，已证的结构定理说明 \(Q_R(M)\) 有限自由，短正合列
\[
0\longrightarrow t_R(M)\longrightarrow M
\longrightarrow Q_R(M)\longrightarrow0
\]
可以分裂。但选择一个分裂比确定挠子模和商模要求更多，它一般不能对全部同态自然地进行。

\begin{proposition}\label{b4:prop:no-natural-torsion-section}
对每个有限生成交换群 \(A\)，记 \(t(A)\) 为其挠子群、\(q_A:A\rightarrow A/t(A)\) 为商映射。在有限生成交换群及群同态的范畴中，不存在一族群同态截面
\[
s_A:A/t(A)\longrightarrow A,\qquad q_As_A=1_{A/t(A)},
\]
使 \(s\) 对每个群同态自然。
\end{proposition}
\begin{proof}
只需考察 \(A=\mathbb Z\oplus\mathbb Z/3\mathbb Z\)。把商识别为 \(\mathbb Z\)，任一截面满足
\(s_A(1)=(1,\bar c)\)。自同构
\[
u:A\longrightarrow A,\qquad u(n,\bar a)=(n,\bar a+\bar n)
\]
在商 \(A/t(A)\) 上诱导恒等。若截面自然，就有 \(us_A=s_A\)。在 \(1\) 处比较得到
\((1,\bar c+\bar1)=(1,\bar c)\)，矛盾。
\end{proof}

对这里讨论的每个有限生成交换群，其挠子群商序列都存在截面；命题否定的是把这些截面选成对全部群同态自然的一族。这也回答了章首提出的问题：各个群分别存在的分解，未必能够同时与所有群同态相容。自然的子对象和自然的商对象，并不自动给出一个自然直和分解。后续的正合性与导出函子需要记录的，正是这些映射和相容关系，而不仅是对象的同构类型。

% !TeX root = ../../Book4.tex
% 纲要标题：### 1.5 范畴代数与图表表示
% 纲要细目（写作范围）：
% ---
% #### 1.5.1 复合的线性化与范畴代数
% #### 1.5.2 图表与左模的对应
% #### 1.5.3 群、路径与偏序的例子

\section{范畴代数与图表表示}
\label{b4:sec:0105}

给出两个 \(k\)-模 \(V_0,V_1\) 及一个线性映射 \(u:V_0\to V_1\)，便得到最简单的模图表。令 \(V=V_0\oplus V_1\)，记投影为 \(p_i:V\rightarrow V_i\)，包含为 \(\iota_i:V_i\rightarrow V\)。作用在总模 \(V\) 上的是幂等自同态 \(e_i=\iota_i\circ p_i\)，它取出第 \(i\) 个分量再放回原位；映射 \(u\) 则给出自同态 \(\iota_1\circ u\circ p_0\)，把第 \(0\) 个分量送入第 \(1\) 个分量，并在第 \(1\) 个分量上为零。这三个自同态之间的关系，已经记录了原来的图表。对一般的有限对象范畴，也可以把全部态射的复合写成一个代数的乘法，使图表成为这个代数上的一个模。

\subsection{复合的线性化与范畴代数}
\label{b4:subsec:0105:01}

\begin{definition}\label{b4:def:category-algebra}
设 \(k\) 为含幺交换环，\(\mathcal C\) 为有有限个对象的小范畴。以 \(\mathcal C\) 的全部态射为基构造自由 \(k\)-模，并在基元素上规定
\[
gf=\begin{cases}
g\circ f,&f:X\to Y,\ g:Y\to Z,\\
0,&f,g\;\text{不能依此次序复合}.
\end{cases}
\]
把乘法双线性地延拓到整个模，所得 \(k\)-代数称为 \(\mathcal C\) 的\term[fanchou-daishu]{范畴代数}[category algebra]，记为 \(k\mathcal C\)。这里每个元素只含有限多个态射的线性组合；并不要求 \(\mathcal C\) 只有有限多个态射。
\end{definition}
\symidx{\(k\mathcal C\)：范畴代数}

\begin{proposition}\label{b4:prop:category-algebra-unit}
设 \(k\) 为含幺交换环，\(\mathcal C\) 为有限对象的小范畴。上述乘法使 \(k\mathcal C\) 成为结合含幺代数，其单位为
\[
1=\sum_{X\in\operatorname{Ob}\mathcal C}e_X,
\qquad e_X=1_X.
\]
各 \(e_X\) 两两正交且幂等。
\end{proposition}
\begin{proof}
对三个基态射 \(f,g,h\)，若它们依次可复合，则 \(h(gf)=(hg)f\) 正是范畴中的结合律。若 \(f,g\) 不能依此次序复合，则 \(f\) 的目标不等于 \(g\) 的源；即使 \(hg\) 可复合，其源仍是 \(g\) 的源，所以 \((hg)f=0\)，两边都为零。若 \(g,h\) 不能依此次序复合，则 \(g\) 的目标不等于 \(h\) 的源；即使 \(gf\) 可复合，其目标仍是 \(g\) 的目标，所以 \(h(gf)=0\)，两边也都为零。双线性延拓遂给出全部元素的结合律。

若 \(f:X\to Y\)，则 \(e_Yf=f=fe_X\)，而其余对象的恒等态射在相应一侧与 \(f\) 相乘均为零。因此所示有限和在两侧都是单位。恒等态射的复合还给出 \(e_X^2=e_X\) 及 \(e_Ye_X=0\)（\(X\ne Y\)）。若范畴为空，所得代数为零代数，单位为 \(0\)。
\end{proof}

有限对象条件用在单位的构造上。若对象无限，则形式和 \(\sum_Xe_X\) 不属于这个只允许有限支集的自由模，不能把它当作代数中的一个元素。另一方面，只有一个对象的群或幺半群，即使有无限多个态射，也仍然给出含幺范畴代数。

\subsection{图表与左模的对应}
\label{b4:subsec:0105:02}

一个协变函子 \(F:\mathcal C\to k\text{-}\mathbf{Mod}\) 给每个对象 \(X\) 配置一个 \(k\)-模 \(F(X)\)，给每个态射 \(f:X\to Y\) 配置一个线性映射 \(F(f):F(X)\to F(Y)\)，并使这些映射服从 \(\mathcal C\) 的复合关系。这样的函子就是一个以 \(\mathcal C\) 为形状的模图表。图表之间的自然变换则是一族与所有图表箭头相容的线性映射。

\begin{theorem}\label{b4:thm:category-algebra-representations}
设 \(k\) 为含幺交换环，\(\mathcal C\) 为有有限个对象的小范畴。以自然变换为态射的函子范畴与含幺左模范畴之间有等价
\[
\operatorname{Fun}(\mathcal C,k\text{-}\mathbf{Mod})
\simeq k\mathcal C\text{-}\mathbf{Mod}.
\]
这个等价把 \(F\) 送到 \(\bigoplus_XF(X)\)，把左 \(k\mathcal C\)-模 \(M\) 送到图表 \(X\mapsto e_XM\)。
\end{theorem}
\begin{proof}
对函子 \(F\)，令 \(M_F=\bigoplus_XF(X)\)。基态射 \(f:X\to Y\) 在 \(X\) 分量上按 \(F(f)\) 作用，所得向量放入 \(Y\) 分量；在其余分量上作用为零。线性延拓定义整个 \(k\mathcal C\) 的作用。若 \(g\) 与 \(f\) 可复合，函子公理给出 \(g(fv)=(gf)v\)；若不可复合，\(fv\) 所在的分量被 \(g\) 送到零。因此作用满足结合律，\(\sum_Xe_X\) 则在各分量上为恒等，故这是含幺左模。

给定自然变换 \(\alpha:F\Rightarrow G\)，映射 \(\bigoplus_X\alpha_X:M_F\to M_G\) 为 \(k\mathcal C\)-线性，因为对每个 \(f:X\to Y\)，自然性恰是
\[
\alpha_YF(f)=G(f)\alpha_X.
\]
直和映射保持恒等与复合，于是得到一个从图表到左模的函子。

反过来，设 \(M\) 为含幺左 \(k\mathcal C\)-模。正交幂等元给出
\[
M=\bigoplus_Xe_XM.
\]
事实上，对象只有有限多个，由 \(1=\sum_Xe_X\) 得到有限和 \(m=\sum_Xe_Xm\)；若 \(\sum_Xm_X=0\) 且 \(m_X\in e_XM\)，乘以 \(e_Y\) 就得到 \(m_Y=0\)。对 \(f:X\to Y\)，定义
\[
F_M(f):e_XM\longrightarrow e_YM,
\qquad m\longmapsto fm.
\]
等式 \(e_Yf=f=fe_X\) 保证值落在指定分量，代数的乘法保证复合相容，\(e_X\) 在 \(e_XM\) 上是恒等。这给出协变函子。模同态 \(h:M\to N\) 满足 \(h(e_Xm)=e_Xh(m)\)，所以可限制到各分量，并与上述箭头相容；因而得到一个从 \(k\mathcal C\) 模范畴回到图表范畴的协变函子。

从 \(F\) 出发再取分量，得到的 \(e_XM_F\) 就是原来的 \(F(X)\)，其箭头与自然变换也相同。从 \(M\) 出发再取直和，则求和映射
\[
\bigoplus_Xe_XM\longrightarrow M,
\qquad (m_X)_X\longmapsto\sum_Xm_X
\]
是模同构，其逆为 \(m\mapsto(e_Xm)_X\)。因对象有限，这族分量自动具有有限支集，确实属于上述直和。这两种识别都与原来的态射相容，故给出拟逆及所需自然同构。
\end{proof}

这里采用左模，是因为态射 \(f:X\to Y\) 应把 \(X\) 分量送入 \(Y\) 分量，而乘法 \(gf\) 表示先作 \(f\) 再作 \(g\)。这一对应的有限对象形式可对照 Webb 的范畴代数构造\cite[Section 2, Proposition 2.1]{Webb2007CategoryRepresentations}。

\subsection{群、路径与偏序的例子}
\label{b4:subsec:0105:03}

把群 \(G\) 看成单对象范畴 \(\mathcal BG\)，便有 \(k\mathcal BG=kG\)：基为群元素，乘法为群乘法。因此，群在 \(k\)-模上的线性表示就是 \(\mathcal BG\) 形状的模图表。自然变换的唯一分量与每个群元素的作用相容，正是表示之间的同态。这里没有要求 \(G\) 有限。
\symidx{\(kG\)：群代数}

\ProseParagraphBreak
再取一个有有限多个顶点的有向图 \(Q\)。以顶点为对象，以有向路径为态射，空路径为恒等，路径的连接为复合，得到自由路径范畴 \(\mathcal P(Q)\)。其范畴代数称为\term[lujing-daishu]{路径代数}[path algebra]，记为 \(kQ\)。从 \(\mathcal P(Q)\) 到 \(k\text{-}\mathbf{Mod}\) 的函子由各顶点的模和各箭头的线性映射唯一决定：任意路径的映射是沿路径依次复合，空路径映成恒等。因而 \(kQ\)-模等价于这个有向图的线性表示。图若有有向环，路径通常有无限多个，但顶点仍有限，所以路径代数仍有单位。
\symidx{\(\mathcal P(Q),\ kQ\)：有向图的路径范畴与路径代数}

\ProseParagraphBreak
对于单箭头 \(0\xrightarrow{a}1\)，基只有 \(e_0,e_1,a\)。对应
\[
e_0\longmapsto E_{11},\qquad
e_1\longmapsto E_{22},\qquad
a\longmapsto E_{21}
\]
把路径代数识别为 \(2\times2\) 下三角矩阵代数。这里 \(E_{ij}\) 表示第 \(i\) 行第 \(j\) 列为 \(1\)、其余为零的矩阵。左模中的两个幂等分量给出 \(V_0,V_1\)，而 \(a\) 给出 \(u:V_0\to V_1\)，恢复节首的图表。若只有一个顶点与一个环箭头 \(a\)，则每条路径唯一写成 \(1,a,a^2,\ldots\)，其中 \(1\) 是空路径。路径乘法满足 \(a^ia^j=a^{i+j}\)，所以 \(a\mapsto t\) 把路径代数识别为 \(k[t]\)；只有这一个生成箭头，并不会出现不同箭头次序的非交换问题。一个左模就是一个 \(k\)-模连同一个指定自同态。

\ProseParagraphBreak
设 \(P\) 为有限偏序集，把它视为偏序范畴。对 \(x\le y\)，记唯一态射为 \(a_{xy}:x\to y\)。于是 \(kP\) 的基由全部这样的区间指标给出，并满足
\[
a_{yz}a_{xy}=a_{xz},
\]
其余不能复合的基元素乘积为零。把偏序线性排列，使 \(x<y\) 时 \(x\) 排在 \(y\) 前，映射 \(a_{xy}\mapsto E_{yx}\) 把 \(kP\) 写成一个下三角矩阵子代数。这是有限偏序集的\term[guanlian-daishu]{关联代数}[incidence algebra]的一种复合约定：矩阵的列指标为源、行指标为目标。若采用把区间 \((x,y)\) 放在第 \(x\) 行、第 \(y\) 列的通常卷积约定，则对应的是此处代数的相反代数。
\symidx{\(kP\)：有限偏序集的关联代数}

\ProseParagraphBreak
偏序图表与自由路径图表之间还有实质区别。在菱形偏序 \(0<1<3\)、\(0<2<3\) 中，从 \(0\) 到 \(3\) 只有一个态射，故两条复合路径必须给出相同线性映射；在只有这四条箭头的自由路径范畴中，两条路径是不同态射，没有这一等式。相应地，前者的范畴代数是后者的路径代数模去“两条路径之差”所得的商。这样，范畴中的复合关系既能约束图表，也能写成代数中的关系。


\begin{chaptersummary}
范畴保留对象之间的态射及复合，不要求每个对象都有预先指定的底集合。单态射和满态射因此用消去性质定义；在模范畴中它们分别对应单射和满射，在环范畴中则不能照搬后一解释。

\ProseParagraphBreak
函子必须说明怎样作用于态射，自然变换必须对源范畴中的全部态射满足方块条件。变换逐分量复合，逐分量取逆也仍然自然。由此可以把函子本身作为对象；源范畴的小性保证自然变换组成集合。

\ProseParagraphBreak
全忠实性控制 Hom 集，本质满控制对象的同构类型。选择代表及同构以后，这些条件给出拟逆；小范畴中的选择由通常选择公理承担，大范畴则须明确指定相应数据。切片与逗号范畴把附加映射纳入对象，为后续泛性质提供统一的叙述方式。

\ProseParagraphBreak
自然性不等于唯一性，也不禁止一切选择。它总是相对于两个确定函子而言。对整环上的模，挠子模及无挠商的构造对模同态自然。在主理想整环上的有限生成情形，相应短正合列分裂，但分裂的存在不保证能够自然地选择分裂；有限生成交换群已经给出了这种失败的例子。分类对象与比较映射因而是不同任务。

\ProseParagraphBreak
对于有限对象的小范畴，恒等态射形成范畴代数的一组正交幂等元。图表的直和给出左模，幂等分量又把左模拆回原图表；自然变换对应模同态。群代数、路径代数与有限偏序集的关联代数都是这种构造的实例，而路径之间的关系对应代数中的关系。
\end{chaptersummary}

\BookExercises

\begin{exercise}\label{b4:ex:01-additive-monoid-category}
把加法幺半群 \((\mathbb N,0,+)\) 看作单对象范畴。求其同构态射、单态射与满态射；再把它与偏序范畴 \((\mathbb N,\le)\) 比较，说明相同的符号集合为什么可以给出不同范畴。
\end{exercise}
\begin{solution}
在单对象范畴中，复合为加法，恒等态射是 \(0\)。只有 \(0\) 有加法逆，所以只有恒等态射可逆。自然数加法左右可消去，因此每个态射都单且满。偏序范畴则有一个对象对应每个自然数，只有当 \(m\le n\) 时才有从 \(m\) 到 \(n\) 的唯一箭头；其同构也只有恒等，所有箭头同样单且满。两者的对象类和 Hom 集完全不同，不能只凭都使用自然数就把它们混同。
\end{solution}

\begin{exercise}\label{b4:ex:01-preorder-quotient}
设 \(P\) 为集合上的预序，定义 \(x\sim y\) 当且仅当 \(x\le y\) 且 \(y\le x\)。证明 \(P/{\sim}\) 是偏序集，并把商映射解释为范畴等价，注明选择所用的位置。
\end{exercise}
\begin{solution}
反身与传递性使 \(\sim\) 为等价关系。定义 \([x]\le[y]\) 当且仅当 \(x\le y\)；更换代表时由两端的双向关系及传递性可知条件不变。反对称性来自 \(\sim\) 的定义。商函子在每对对象的 Hom 上都是两个空集之间或两个单点集之间的双射，且每个等价类有代表，所以全忠实且本质满。对象集为集合，选择公理可为每类取一个代表，再由等价判据构造拟逆。选择只用于选代表，不用于商偏序自身的定义。
\end{solution}

\begin{exercise}\label{b4:ex:01-group-functors-intertwiners}
设 \(G,H\) 为群。把函子 \(\mathcal BG\rightarrow\mathcal BH\) 与群同态对应。若函子由 \(\rho,\sigma:G\rightarrow H\) 给出，描述全部自然变换 \(\rho\Rightarrow\sigma\)，并判断何时二者自然同构。若 \(G\) 由子集 \(T\) 生成，证明只需对 \(g\in T\) 检验自然性；再用一个例子说明，在任意子集上检验通常不够。
\end{exercise}
\begin{solution}
对象映射没有选择，态射映射保持复合与恒等，恰好是保单位的乘法同态，也就是群同态。自然变换只有一个分量 \(h\in H\)，对每个 \(g\in G\) 的自然性是
\[
\sigma(g)h=h\rho(g).
\]
因此它等价于 \(\sigma(g)=h\rho(g)h^{-1}\) 对所有 \(g\) 成立。因为 \(H\) 中每个元素可逆，每个自然变换都是自然同构；存在自然同构恰好表示两个同态在 \(H\) 中共轭。自变换的分量构成 \(\rho(G)\) 在 \(H\) 中的中心化子群。

映射 \(g\mapsto h\rho(g)h^{-1}\) 与 \(\sigma\) 都是群同态。在生成元上相等便在其逆元及有限乘积上相等，故只需检验 \(T\)。若只在任意子集上检验，则取 \(G=(\mathbb Z/2\mathbb Z)^2\)、\(H=\mathbb Z/2\mathbb Z\)，令 \(\rho,\sigma\) 为两个坐标投影。它们在对角子群上相等，却不在整个 \(G\) 上相等。由于 \(H\) 交换，共轭不改变同态，所以不存在自然同构。
\end{solution}

\begin{exercise}\label{b4:ex:01-powerset-variance}
证明幂集赋值 \(X\mapsto\mathcal P(X)\) 可以通过直接像成为协变函子，也可以通过逆像成为反变函子。证明单元素子集给出
\(\operatorname{Id}_{\mathbf{Set}}\Rightarrow\mathcal P\) 的自然变换，其中 \(\mathcal P\) 取直接像；再证明取补集是逆像版本的自然自同构。
\end{exercise}
\begin{solution}
对函数 \(f:X\rightarrow Y\)，直接像为 \(A\mapsto f(A)\)，逆像为
\(B\mapsto f^{-1}(B)\)。等式
\((gf)(A)=g(f(A))\) 与
\((gf)^{-1}(C)=f^{-1}(g^{-1}(C))\)
分别给出协变与反变复合，恒等情况直接成立。因 \(f(\{x\})=\{f(x)\}\)，单元素子集的方块交换。对补集映射 \(c_X(A)=X\setminus A\)，有
\[
f^{-1}(Y\setminus B)=X\setminus f^{-1}(B),
\]
所以它在反变版本中自然，且 \(c_X^2=1\)。它对直接像一般不自然。例如对空函数 \(f:\varnothing\rightarrow\{*\}\)，记直接像为 \(f_*\)，两条复合路径在 \(\varnothing\in\mathcal P(\varnothing)\) 处分别给出
\[
c_{\{*\}}\bigl(f_*(\varnothing)\bigr)=\{*\},
\qquad
f_*\bigl(c_{\varnothing}(\varnothing)\bigr)=f_*(\varnothing)=\varnothing.
\]
这里在 \(\varnothing\) 内取补集仍得 \(\varnothing\)，而在 \(\{*\}\) 内取其补集得 \(\{*\}\)，故自然性方块不交换。
\end{solution}

\begin{exercise}\label{b4:ex:01-abelianization-functor}
对群 \(G\)，记 \([G,G]\) 为交换子子群。证明 \(G\mapsto[G,G]\) 可以由限制同态成为群范畴上的函子，且交换化 \(G\mapsto G/[G,G]\) 成为取值于交换群的函子。说明它与\subsecref{b4:subsec:0102:03}末尾的群中心构造有何不同。
\end{exercise}
\begin{solution}
群同态 \(f:G\rightarrow H\) 满足
\(f([x,y])=[f(x),f(y)]\)，所以把生成交换子子群的元素送入 \([H,H]\)，进而限制同态良定义。恒等和复合与限制相容。它也诱导
\(\bar f:G/[G,G]\rightarrow H/[H,H]\)，商群交换；诱导映射保持恒等与复合，所以给出函子，商映射满足
\(\bar f q_G=q_Hf\)。中心的定义要求与目标群中所有元素交换，而一般同态只控制像中的元素；交换子生成元的像则直接仍是交换子，因此后者的函子规则适用于全部同态。
\end{solution}

\begin{exercise}\label{b4:ex:01-n-torsion-quotient}
固定整数 \(n\ge2\)。对交换群 \(A\)，令 \(A[n]=\{a:na=0\}\)。
证明包含后取商给出自然变换
\[
A[n]\longrightarrow A/nA.
\]
用 \(A=\mathbb Z\) 与 \(A=\mathbb Z/n^2\mathbb Z\) 说明其分量不必满，也不必单。
\end{exercise}
\begin{solution}
同态把被 \(n\) 零化的元素送到同样的元素，也把 \(nA\) 送入 \(nB\)，故两端均为函子，映射
\(a\mapsto a+nA\) 良定义。对 \(f:A\rightarrow B\)，两条路径都把 \(a\) 送到 \(f(a)+nB\)，所以自然。若 \(A=\mathbb Z\)，源为零、目标为 \(\mathbb Z/n\mathbb Z\)，不满；若 \(A=\mathbb Z/n^2\mathbb Z\)，源是非零子群 \(nA\)，它在目标中全部变成零，所以不单。
\end{solution}

\begin{exercise}\label{b4:ex:01-poset-natural-transformations}
设 \(P,Q\) 为偏序集，\(F,G:P\rightarrow Q\) 为保序映射。证明存在自然变换 \(F\Rightarrow G\) 当且仅当每个 \(p\in P\) 都有 \(F(p)\le G(p)\)，并证明这样的变换唯一。何时它是自然同构？
\end{exercise}
\begin{solution}
在 \(p\) 处有分量箭头，恰好等价于 \(F(p)\le G(p)\)。若每处分量存在，所有自然性方块的两条复合路径都有相同起点和终点；偏序范畴中这样的箭头至多一个，所以方块必交换。分量本来就唯一，因此变换唯一。分量可逆等价于 \(F(p)=G(p)\)，故自然同构存在当且仅当两个保序映射逐点相等。
\end{solution}

\begin{exercise}\label{b4:ex:01-interchange-law}
设 \(F_0,F_1,F_2:\mathcal C\rightarrow\mathcal D\)、
\(H_0,H_1,H_2:\mathcal D\rightarrow\mathcal E\)，以及
\[
\eta:F_0\Rightarrow F_1,\quad\theta:F_1\Rightarrow F_2,\quad
\alpha:H_0\Rightarrow H_1,\quad\beta:H_1\Rightarrow H_2.
\]
证明水平复合与垂直复合满足
\((\beta\alpha)*(\theta\eta)=(\beta*\theta)(\alpha*\eta)\)。
\end{exercise}
\begin{solution}
两端都从 \(H_0F_0\) 到 \(H_2F_2\)。右端在 \(X\) 处为
\[
\beta_{F_2X}H_1(\theta_X)\alpha_{F_1X}H_0(\eta_X).
\]
对箭头 \(\theta_X:F_1X\rightarrow F_2X\)，\(\alpha\) 的自然性给出
\(H_1(\theta_X)\alpha_{F_1X}=\alpha_{F_2X}H_0(\theta_X)\)。
代入后得到
\[
\beta_{F_2X}\alpha_{F_2X}H_0(\theta_X\eta_X),
\]
这正是左端的分量。每个对象处分量相同，故两自然变换相等。
\end{solution}

\begin{exercise}\label{b4:ex:01-three-conditions-independent}
分别给出满足以下两项而不满足第三项的有限范畴间函子：
全且本质满但不忠实；忠实且本质满但不全；全忠实但不本质满。
\end{exercise}
\begin{solution}
第一例取 \(\mathcal B(\mathbb Z/2\mathbb Z)\rightarrow\mathbf1\)，它在唯一 Hom 上满射且对象本质满，却把两个不同态射合并。第二例取两个对象的离散范畴到两元素链 \(0<1\) 的恒等对象映射；每个 Hom 上单射且对象全被包括，但源中没有提升 \(0\rightarrow1\) 的箭头。第三例取单对象离散范畴到两对象离散范畴的包含；保留对象之间的 Hom 全相同，但另一个对象没有同构原像。
\end{solution}

\begin{exercise}\label{b4:ex:01-skeletal-equivalence}
若一个范畴中同构对象必相等，称它为骨架范畴。证明两个小骨架范畴之间的范畴等价其实是范畴同构。
\end{exercise}
\begin{solution}
设等价为 \(F:\mathcal C\rightarrow\mathcal D\)。本质满和 \(\mathcal D\) 的骨架性说明每个 \(Y\) 实际等于某个 \(FX\)，所以对象映射满射。若 \(FX=FX'\)，全忠实反射同构，得 \(X\cong X'\)，再由 \(\mathcal C\) 的骨架性得 \(X=X'\)，故对象映射双射。对对象采用真正逆函数，对态射采用各 Hom 双射的逆，函子性由 \(F\) 保持恒等与复合及单射性推出。这个逆与 \(F\) 的两次复合严格为恒等。
\end{solution}

\begin{exercise}\label{b4:ex:01-composite-equivalence}
设 \(F:\mathcal C\rightarrow\mathcal D\)、\(H:\mathcal D\rightarrow\mathcal E\) 分别有拟逆 \(G,K\)。给定自然同构
\(\eta:\operatorname{Id}_{\mathcal C}\Rightarrow GF\)、\(\varepsilon:FG\Rightarrow\operatorname{Id}_{\mathcal D}\)、
\(\mu:\operatorname{Id}_{\mathcal D}\Rightarrow KH\)、\(\nu:HK\Rightarrow\operatorname{Id}_{\mathcal E}\)。
直接写出 \(HF\) 的拟逆与两族自然同构。
\end{exercise}
\begin{solution}
拟逆取 \(GK\)。从 \(\operatorname{Id}_{\mathcal C}\) 到 \(GKHF\) 的自然同构为
\[
X\xrightarrow{\eta_X}GFX
\xrightarrow{G(\mu_{FX})}GKHFX.
\]
从 \(HFGK\) 到 \(\operatorname{Id}_{\mathcal E}\) 的自然同构为
\[
HFGKZ\xrightarrow{H(\varepsilon_{KZ})}HKZ
\xrightarrow{\nu_Z}Z.
\]
每个分量为同构，且两族分别由自然变换与函子的左右复合及逐分量复合得到，所以自然。这直接验证范畴等价可复合。
\end{solution}

\begin{exercise}\label{b4:ex:01-arrow-poset-count}
令 \([2]\) 为链 \(0<1<2\) 对应的范畴。描述函子范畴 \(\operatorname{Fun}([1],[2])\)，列出其对象，并计算全部态射的数目。
\end{exercise}
\begin{solution}
对象为有序对 \((i,j)\)，其中 \(0\le i\le j\le2\)，共六个：
\[
(0,0),(0,1),(0,2),(1,1),(1,2),(2,2).
\]
由偏序自然变换的判据，\((i,j)\) 到 \((i',j')\) 有唯一态射，当且仅当 \(i\le i'\)、\(j\le j'\)。所以这个函子范畴本身为逐坐标比较的偏序范畴。按所列顺序，各对象发出的态射数分别为 \(6,5,3,3,2,1\)，总数为 \(20\)，其中包括六个恒等态射。
\end{solution}

\begin{exercise}\label{b4:ex:01-slice-two-points}
设 \(S=\{0,1\}\)。构造明确的等价
\(\mathbf{Set}/S\simeq\mathbf{Set}\times\mathbf{Set}\)，不得通过为全部集合选择代表来构造拟逆。
\end{exercise}
\begin{solution}
把 \(p:X\rightarrow S\) 送到两条纤维
\((p^{-1}(0),p^{-1}(1))\)，切片态射因保持 \(p\) 而分别限制到两纤维。这给出一个函子。反向把 \((A,B)\) 送到带标记不交并
\((\{0\}\times A)\sqcup(\{1\}\times B)\rightarrow S\)，映射按标记所在分量作用。这是具体定义的函子，不用选择。先取纤维再不交并，有自然双射
\((p(x),x)\mapsto x\)；反向复合得到的纤维由去掉固定标记自然同构于 \(A,B\)。两族映射与各分量函数相容，所以给出等价。
\end{solution}

\begin{exercise}\label{b4:ex:01-actions-as-functors}
设 \(G\) 为群。证明函子 \(\mathcal BG\rightarrow\mathbf{Set}\) 等价于一个左 \(G\)-集合结构，自然变换等价于等变映射。
\end{exercise}
\begin{solution}
函子给唯一对象指定集合 \(X\)，给每个 \(g\in G\) 指定函数 \(\rho(g):X\rightarrow X\)，满足
\(\rho(gh)=\rho(g)\rho(h)\)、\(\rho(1)=1_X\)。令 \(gx=\rho(g)(x)\)，这就是左作用；每个 \(\rho(g)\) 的逆为 \(\rho(g^{-1})\)。反过来，左作用的两条公理给出函子。两个这样的函子间自然变换只有一个分量 \(u:X\rightarrow Y\)，自然性为
\(u(gx)=g\,u(x)\)，恰是等变性。
\end{solution}

\begin{exercise}\label{b4:ex:01-comma-tautological-map}
设 \(S:\mathcal A\to\mathcal C\)、\(T:\mathcal B\to\mathcal C\) 为函子。对逗号范畴 \((S\downarrow T)\) 的投影
\(P:(S\downarrow T)\rightarrow\mathcal A\)、
\(Q:(S\downarrow T)\rightarrow\mathcal B\)，证明对象中的第三分量给出自然变换 \(SP\Rightarrow TQ\)。说明在箭头范畴中这族分量是什么。
\end{exercise}
\begin{solution}
在对象 \((A,B,u)\) 处分量就是 \(u:S(A)\rightarrow T(B)\)。
对逗号态射 \((a,b)\)，自然性恰为定义中的等式
\(T(b)u=u'S(a)\)，故成立。在箭头范畴中，\(P,Q\) 是取起点与终点的函子，分量就是被研究的箭头自身；交换方块正是这一变换的自然性。
\end{solution}

\begin{exercise}\label{b4:ex:01-natural-subset-choices}
设 \(\mathcal P:\mathbf{Set}\rightarrow\mathbf{Set}\) 使用直接像。求所有自然变换
\(\eta:\operatorname{Id}_{\mathbf{Set}}\Rightarrow\mathcal P\)。
\end{exercise}
\begin{solution}
在单点集 \(\{*\}\) 上，\(\eta_{\{*\}}(*)\) 只有两个可能值：空集或整个单点集。对任意 \(x\in X\)，用函数 \(a_x:\{*\}\rightarrow X\) 的自然性，得
\[
\eta_X(x)=\mathcal P(a_x)\bigl(\eta_{\{*\}}(*)\bigr).
\]
这里 \(\mathcal P(a_x)\) 是取子集直接像的映射；\(a_x\) 本身的输入则是单点集中的元素。因此第一种情形所有分量把元素送到空集，第二种情形把 \(x\) 送到 \(\{x\}\)。这两族对任意函数都自然，空集处的函数也唯一，所以恰有两种自然变换。
\end{solution}

\begin{exercise}\label{b4:ex:01-frobenius-natural}
固定素数 \(p\)，考虑交换 \(\mathbb F_p\)-代数及保幺 \(\mathbb F_p\)-代数同态的范畴。证明
\(a\mapsto a^p\) 给出恒等函子的自然自变换，并用一个例子说明它不必是自然同构。
\end{exercise}
\begin{solution}
交换性与二项式系数模 \(p\) 消失使
\((a+b)^p=a^p+b^p\)，并有 \((ab)^p=a^pb^p\)、\(1^p=1\)。
对 \(c\in\mathbb F_p\)，\(c^p=c\)，所以这还是 \(\mathbb F_p\)-代数同态。任意同态 \(f\) 满足 \(f(a^p)=f(a)^p\)，故自然。在 \(\mathbb F_p[t]\) 上，像只有次数为 \(p\) 的倍数的单项式的组合，\(t\) 不在像中，因此这一分量不满，不是自然同构。
\end{solution}

\begin{exercise}\label{b4:ex:01-rationalization-not-reflect}
使用第二卷第6.3.3小节的扩张标量函子（取环同态 \(\mathbb Z\hookrightarrow\mathbb Q\)）
\(\mathbb Q\otimes_{\mathbb Z}-:\mathbf{Ab}\rightarrow\mathbf{Vect}_{\mathbb Q}\)，给出一个不是同构但其像为同构的群同态。说明为什么这不与全忠实反射同构的命题冲突。
\end{exercise}
\begin{solution}
取 \(f:\mathbb Z\rightarrow\mathbb Z\)、\(f(n)=2n\)。它不满，因而不是同构；张量后在
\(\mathbb Q\otimes\mathbb Z\simeq\mathbb Q\) 上为乘 \(2\)，逆为乘 \(1/2\)。此函子不全忠实：例如 \(\mathbb Z/2\mathbb Z\) 张量后为零，其恒等同态与零同态有相同的像，故甚至不忠实。保留同构适用于任意函子，反射同构则需要额外条件。
\end{solution}

\begin{exercise}\label{b4:ex:01-primary-decomposition-natural}
对有限交换群 \(A\) 和素数 \(p\)，令
\(A_{(p)}=\{a:\;\text{某个}\;p^r\;\text{零化}\;a\}\)。
% 跨卷引用目标：b1:thm:abelian-primary（9.2.4）；b2:thm:pid-elementary-divisors（4.3.4）。
利用第一卷定理9.2.4的有限交换群主分解或第二卷定理4.3.4的 PID 模初等因子分解，说明自然映射
\(\bigoplus_pA_{(p)}\rightarrow A\) 是自然同构。解释为什么这与\cref{b4:prop:no-natural-torsion-section}中不存在自然挠子群商映射截面的结论不同。
\end{exercise}
\begin{solution}
同态保留被 \(p\) 的幂零化的元素，因此各 \(A_{(p)}\) 是函子，其包含自然。已证的初等因子分解说明每个有限交换群是有限多个素数幂阶循环群的直和；在该分解下，\(A_{(p)}\) 正是所有 \(p\)-幂阶分量之和。因而只有整除这些循环群阶数之积 \(|A|\) 的素数可能给出非零分量，所示直和实际只有有限多个非零项，各包含之和为双射，也就是同构。映射本身由包含定义，不依赖用于证明双射的分类同构；它对所有同态自然。这里各主分量由元素的零化性质内在确定，正文的自由补群则没有这种自然指定。
\end{solution}

\begin{exercise}\label{b4:ex:01-two-coordinate-functors}
设 \(k\) 为域，\(\mathcal V\) 为有限维 \(k\)-向量空间范畴的一个小全子范畴。对每个 \(V\in\mathcal V\)，分别固定两组坐标同构
\(\theta_V:V\rightarrow k^{\dim V}\)、\(\psi_V:V\rightarrow k^{\dim V}\)。定义函子 \(F_\theta,F_\psi:\mathcal V\to\mathbf{Vect}_k^{\mathrm{fd}}\)：两者均把 \(V\) 送到 \(k^{\dim V}\)，并分别把 \(f:V\to W\) 送到 \(\theta_Wf\theta_V^{-1}\) 与 \(\psi_Wf\psi_V^{-1}\)。写出它们之间的自然同构，计算每个线性映射的两个矩阵之间的关系。
\end{exercise}
\begin{solution}
取 \(P_V=\psi_V\theta_V^{-1}\)。对 \(f:V\rightarrow W\)，有
\[
F_\psi(f)=\psi_Wf\psi_V^{-1}
=P_WF_\theta(f)P_V^{-1}.
\]
因此 \(F_\psi(f)P_V=P_WF_\theta(f)\)，正是
\(P:F_\theta\Rightarrow F_\psi\) 的自然性。每个 \(P_V\) 可逆，故为自然同构。若 \(V=W\) 且只考察一个算子，等式退化为熟悉的相似变换；一般映射的两端空间不同，应保留两个不同的换基矩阵。
\end{solution}

\begin{exercise}\label{b4:ex:01-chain-category-algebra}
设 \(k\) 为域，\(P=\{0<1<2\}\) 为三元素全序集。把 \(kP\) 具体识别为矩阵代数，说明一个左 \(kP\)-模为何等价于两个可复合线性映射 \(V_0\xrightarrow{u}V_1\xrightarrow{v}V_2\)。写出模同态对应的三个线性映射及其条件。
\end{exercise}
\begin{solution}
基态射为 \(e_0,e_1,e_2,a_{01},a_{12},a_{02}\)。映射 \(e_i\mapsto E_{i+1,i+1}\)、\(a_{ij}\mapsto E_{j+1,i+1}\) 给出到全部 \(3\times3\) 下三角矩阵的基双射，且 \(E_{32}E_{21}=E_{31}\) 对应 \(a_{12}a_{01}=a_{02}\)，其余乘积也与复合规则相同，故是代数同构。

左模 \(M\) 分解为 \(V_i=e_iM\)。\(a_{01}\)、\(a_{12}\) 分别给出 \(u,v\)，\(a_{02}\) 的作用被乘法关系迫使等于 \(vu\)；没有额外数据。反向，把这三个空间直和，并按 \(u,v,vu\) 定义各基态射的作用，就得到左模。

模同态必须保留幂等分量，故由 \(h_i:V_i\to W_i\) 给出。若目标箭头为 \(u',v'\)，其线性条件为 \(h_1u=u'h_0\)、\(h_2v=v'h_1\)。它们自动推出 \(h_2vu=v'u'h_0\)，恰好是整张图表的自然性条件。
\end{solution}

\begin{exercise}\label{b4:ex:01-parallel-arrows-algebra}
设 \(k\) 为域，\(Q\) 只有两个顶点 \(0,1\) 和两条箭头 \(a,b:0\to1\)。描述 \(kQ\) 的基与乘法。对 \(V_0=V_1=k\) 的图表，按两个箭头的标量 \((s,t)\in k^2\) 分类其同构类型。
\end{exercise}
\begin{solution}
基为 \(e_0,e_1,a,b\)，单位为 \(e_0+e_1\)。非零的含箭头乘积只有 \(e_1a=a=ae_0\)、\(e_1b=b=be_0\)，任意两个箭头相乘均为零。尤其没有把 \(a,b\) 识别的关系。

两图表 \((s,t)\)、\((s',t')\) 同构，等价于存在两个非零标量 \(c_0,c_1\) 使 \(c_1s=s'c_0\)、\(c_1t=t'c_0\)，即 \((s',t')=c(s,t)\)，其中 \(c=c_1/c_0\in k^\times\)。两个顶点上的换基通过同一个比值 \(c_1/c_0\) 同时缩放两条箭头，不能分别调整 \(s,t\)。所以零对单独成一类，非零对按共同的非零标量倍分类。若 \(s\ne0\)，取共同倍数 \(s^{-1}\)，得到唯一代表 \((1,t/s)\)；若 \(s=0\)，则 \(t\ne0\)，得到代表 \((0,1)\)。非零类保留的是两个箭头之间的比例，因而有一个比例参数以及 \((0,1)\) 这一额外类型；这份数据只看顶点维数是看不到的。
\end{solution}

\begin{exercise}\label{b4:ex:01-diamond-incidence-relation}
设 \(k\) 为域，\(Q\) 的四个顶点与四条箭头为
\[
0\xrightarrow{a}1\xrightarrow{b}3,
\qquad0\xrightarrow{c}2\xrightarrow{d}3.
\]
令 \(P\) 为相应菱形偏序集。证明 \(kP\cong kQ/(ba-dc)\)，其中分母表示所示元素生成的双边理想；分别计算两代数的维数，并描述该关系对图表的要求。
\end{exercise}
\begin{solution}
\(kQ\) 的基包括四条空路径、四条长度一的路径和两条长度二的路径 \(ba,dc\)，没有更长路径，故维数为 \(10\)。在偏序范畴中四条箭头的像如原图，而 \(ba,dc\) 都映成唯一的 \(0\to3\) 态射。这给出满代数同态 \(kQ\to kP\)。\(P\) 中除四个恒等与四条相邻态射外，还有一个 \(0\to3\)，故维数为 \(9\)。上述映射的核是 \(k(ba-dc)\)；这个一维子空间是双边理想，因为除两端恒等态射外，没有可与这两条最长路径复合的非恒等路径。因此核恰为所列双边理想。

若四条箭头的线性映射仍记为 \(a,b,c,d\)，商代数模要求 \(ba=dc\)，即两条从 \(V_0\) 到 \(V_3\) 的复合相等。自由路径图表不要求这一等式，偏序图表则必须满足它。
\end{solution}

\begin{exercise}\label{b4:ex:01-involution-characteristic-two}
设 \(k\) 为特征 \(2\) 的域，\(C_2=\{1,g\}\) 为二阶循环群。证明 \(kC_2\cong k[t]/(t^2)\)，并分类其二维左模的同构类型。
\end{exercise}
\begin{solution}
在群代数中令 \(t=g-1\)，则 \(t^2=g^2-2g+1=0\)。\(1,t\) 是基，故所给商环映射为代数同构。一个二维左模就是一个二维向量空间 \(V\) 连同 \(N:V\to V\)，满足 \(N^2=0\)；\(g\) 作用为 \(1+N\)。

若 \(N=0\)，则任取基后 \(g\) 的矩阵为单位矩阵。若 \(N\ne0\)，选 \(v\) 使 \(u=Nv\ne0\)。\(Nu=0\)，而 \(u,v\) 线性无关：否则 \(v\) 为 \(u\) 的倍数会导致 \(Nv=0\)。它们因而组成基，在此基下
\[
N=\begin{pmatrix}0&1\\0&0\end{pmatrix},
\qquad g=\begin{pmatrix}1&1\\0&1\end{pmatrix}.
\]
这两种类型因 \(N\) 的秩不同而不同构，又已涵盖全部可能，所以二维左模恰有这两个同构类型。
\end{solution}
